\section{O teorema fundamental dos morfismos próprios. Aplicações}
\label{section:III.4}


\subsection{O teorema fundamental}
\label{subsection:III.4.1}

\begin{env}[4.1.1]
\label{III.4.1.1}
Sejam $X$ e $Y$ esquemas ordinários noetherianos, seja $f:X\to Y$ um
morfismo próprio, seja $Y'$ um subconjunto fechado de $Y$, e ponhamos
\[
  X'=f^{-1}(Y').
\]
Denotamos por $\widehat{X}$ e $\widehat{Y}$ os esquemas formais
$X_{/X'}$ e $Y_{/Y'}$, obtidos completando $X$ e $Y$ ao longo de
$X'$ e $Y'$, respectivamente \sref[I]{I.10.8.5}.
Denotamos por $\widehat{f}:\widehat{X}\to\widehat{Y}$ a extensão de
$f$ a estes completamentos \sref[I]{I.10.9.1}.
Para todo $\sh{O}_X$-módulo coerente $\sh{F}$, denotamos por

\oldpage[III]{123}

\[
  \widehat{\sh{F}}=\sh{F}_{/X'}
\]
seu completamento ao longo de $X'$ \sref[I]{I.10.8.4}; este é um
$\sh{O}_{\widehat{X}}$-módulo coerente \sref[I]{I.10.8.8}.
\end{env}

\begin{env}[4.1.2]
\label{III.4.1.2}
Seja $\sh{J}$ um ideal coerente de $\sh{O}_Y$ tal que
\[
  \Supp(\sh{O}_Y/\sh{J})=Y'
\]
\sref[I]{I.5.2.1}.
Sabemos \sref[I]{I.4.4.5} que
\[
  \sh{K}=f^*(\sh{J})\sh{O}_X
\]
é um ideal coerente de $\sh{O}_X$ que satisfaz
\[
  \Supp(\sh{O}_X/\sh{K})=X'.
\]
Para todo $k\geq0$, consideremos o $\sh{O}_X$-módulo coerente
\[
  \sh{F}_k
  =
  \sh{F}\otimes_{\sh{O}_Y}(\sh{O}_Y/\sh{J}^{k+1})
  =
  \sh{F}/\sh{K}^{k+1}\sh{F}.
\]
Os $\sh{O}_Y$-módulos
\[
  \RR^n f_*(\sh{F})
  \quad\text{e}\quad
  \RR^n f_*(\sh{F}_k)
\]
são coerentes para todo $n$ \sref{III.3.2.1}.
Para todo $k\geq0$ e todo $n$, o homomorfismo canônico
$\sh{F}\to\sh{F}_k$ induz, por funtorialidade, um homomorfismo
\[
  \RR^n f_*(\sh{F})
  \longrightarrow
  \RR^n f_*(\sh{F}_k).
  \tag{4.1.2.1}
  \label{III.4.1.2.1}
\]
Como $\sh{F}_k$ é um $\sh{O}_X/\sh{K}^{k+1}$-módulo,
$\RR^n f_*(\sh{F}_k)$ é um
$\sh{O}_Y/\sh{J}^{k+1}$-módulo \sref[0]{0.12.2.1}.
Deduzimos, pois, de \eqref{III.4.1.2.1} um homomorfismo
\[
  \RR^n f_*(\sh{F})
  \otimes_{\sh{O}_Y}(\sh{O}_Y/\sh{J}^{k+1})
  \longrightarrow
  \RR^n f_*(\sh{F}_k).
  \tag{4.1.2.2}
  \label{III.4.1.2.2}
\]
Os dois membros de \eqref{III.4.1.2.2} formam sistemas projetivos.
O limite projetivo do membro esquerdo é o completamento
\[
  \bigl(\RR^n f_*(\sh{F})\bigr)_{/Y'},
\]
que denotaremos por
\[
  \widehat{\RR^n f_*(\sh{F})}.
\]
Os homomorfismos \eqref{III.4.1.2.2} são compatíveis com os sistemas
projetivos e definem, portanto, ao passar ao limite, um homomorfismo
canônico
\[
  \vphi_n:
  \widehat{\RR^n f_*(\sh{F})}
  \longrightarrow
  \varprojlim_k\RR^n f_*(\sh{F}_k).
  \tag{4.1.2.3}
  \label{III.4.1.2.3}
\]
Além disso, \eqref{III.4.1.2.2} é um homomorfismo de
$\sh{O}_Y/\sh{J}^{k+1}$-módulos.
Por \sref[I]{I.10.8.3}, pode, portanto, ser considerado um
homomorfismo contínuo de $\sh{O}_{\widehat{Y}}$-módulos topológicos
munidos de suas topologias pseudodiscretas \sref[0]{0.3.8.1}.
Consequentemente, $\vphi_n$ é um homomorfismo contínuo de
$\sh{O}_{\widehat{Y}}$-módulos topológicos.
\end{env}

\begin{env}[4.1.3]
\label{III.4.1.3}
Seja $i:\widehat{X}\to X$ o morfismo canônico de espaços anelados
definido em \sref[I]{I.10.8.7}.
Temos o diagrama comutativo
\[
  \xymatrix{
    X_k\ar[r]^{h_k}\ar[dr]_{i_k} &
    \widehat{X}\ar[d]^{i}\\
    & X
  }
  \tag{4.1.3.1}
  \label{III.4.1.3.1}
\]
onde $X_k$ é o subesquema fechado de $X$ definido por
$\sh{K}^{k+1}$, $i_k$ é a imersão canônica, e $h_k$ é o
morfismo de espaços anelados que corresponde à identidade sobre
os espaços subjacentes e ao homomorfismo canônico
\[
  \sh{O}_{\widehat{X}}
  \longrightarrow
  \sh{O}_X/\sh{K}^{k+1}
\]
\sref[I]{I.10.5.2}.
Além disso,
\[
  \widehat{\sh{F}}=i^*(\sh{F})
\]
salvo isomorfismo canônico \sref[I]{I.10.8.8}.
Temos
\[
  \HH^n(X_k,i_k^*(\sh{F}_k))
  =
  \HH^n(X,\sh{F}_k)
  \tag{4.1.3.2}
  \label{III.4.1.3.2}
\]
salvo isomorfismo canônico.

\oldpage[III]{124}

Com efeito,
\[
  \sh{F}_k=(i_k)_*(i_k^*(\sh{F}_k))
\]
(G,~II,~4.9.1).
O homomorfismo canônico
\[
  \HH^n(\widehat{X},\widehat{\sh{F}})
  \longrightarrow
  \HH^n(X_k,h_k^*(\widehat{\sh{F}}))
\]
\sref[0]{0.12.1.3.5} pode escrever-se também
\[
  \HH^n(\widehat{X},\widehat{\sh{F}})
  \longrightarrow
  \HH^n(X,\sh{F}_k).
  \tag{4.1.3.3}
  \label{III.4.1.3.3}
\]
Estes homomorfismos formam um sistema projetivo e definem, portanto,
um homomorfismo canônico
\[
  \psi_{n,X}:
  \HH^n(\widehat{X},\widehat{\sh{F}})
  \longrightarrow
  \varprojlim_k\HH^n(X,\sh{F}_k).
  \tag{4.1.3.4}
  \label{III.4.1.3.4}
\]
Substituindo $X$ por um subconjunto aberto da forma $f^{-1}(V)$,
onde $V$ é um aberto afim de $Y$, e tendo em conta
\sref{III.1.4.11}, obtemos homomorfismos
\[
  \psi_{n,V}:
  \HH^n(\widehat{X}\cap f^{-1}(V),\widehat{\sh{F}})
  \longrightarrow
  \varprojlim_k\Gamma(V,\RR^n f_*(\sh{F}_k)).
  \tag{4.1.3.5}
  \label{III.4.1.3.5}
\]
Estes homomorfismos comutam manifestamente com a restrição de $V$
a um subconjunto aberto afim menor.
Definem, portanto, um homomorfismo canônico de feixes
\[
  \psi_n:
  \RR^n\widehat{f}_*(\widehat{\sh{F}})
  \longrightarrow
  \varprojlim_k\RR^n f_*(\sh{F}_k).
  \tag{4.1.3.6}
  \label{III.4.1.3.6}
\]
\end{env}

\begin{env}[4.1.4]
\label{III.4.1.4}
Por fim, seja $j:\widehat{Y}\to Y$ o morfismo canônico de espaços
anelados \sref[I]{I.10.8.7}.
Como $\RR^n f_*(\sh{F})$ é um $\sh{O}_Y$-módulo coerente
\sref{III.3.2.1}, temos, salvo isomorfismo canônico,
\[
  j^*(\RR^n f_*(\sh{F}))
  =
  \widehat{\RR^n f_*(\sh{F})}
\]
\sref[I]{I.10.8.8}.
Existe, consequentemente, um homomorfismo canônico
\[
  \rho_n:
  \widehat{\RR^n f_*(\sh{F})}
  =
  j^*(\RR^n f_*(\sh{F}))
  \longrightarrow
  \RR^n\widehat{f}_*(i^*(\sh{F}))
  =
  \RR^n\widehat{f}_*(\widehat{\sh{F}})
  \tag{4.1.4.1}
  \label{III.4.1.4.1}
\]
definido para os espaços anelados em geral; veja-se a demonstração
de \sref{III.1.4.15}.
Demonstraremos que o diagrama
\[
  \xymatrix{
    \widehat{\RR^n f_*(\sh{F})}
      \ar[rr]^{\rho_n}\ar[dr]_{\vphi_n}
    &&
    \RR^n\widehat{f}_*(\widehat{\sh{F}})
      \ar[dl]^{\psi_n}\\
    &
    \displaystyle\varprojlim_k\RR^n f_*(\sh{F}_k)
    &
  }
  \tag{4.1.4.2}
  \label{III.4.1.4.2}
\]
é comutativo.
Basta demonstrar a comutatividade do diagrama formado pelos
homomorfismos correspondentes de pré-feixes.
Podemos, portanto, restringir-nos ao caso em que $Y$ é afim, e basta
demonstrar que
\[
  \xymatrix{
    \widehat{\HH^n(X,\sh{F})}
      \ar[rr]^{\rho_n}\ar[dr]_{\vphi_n}
    &&
    \HH^n(\widehat{X},\widehat{\sh{F}})
      \ar[dl]^{\psi_{n,X}}\\
    &
    \displaystyle\varprojlim_k\HH^n(X,\sh{F}_k)
    &
  }
  \tag{4.1.4.3}
  \label{III.4.1.4.3}
\]
é comutativo.
A comutatividade de \eqref{III.4.1.3.1}, junto com as relações
entre os grupos de cohomologia estabelecidas em \sref{III.4.1.3}, dá
o diagrama comutativo

\oldpage[III]{125}

\[
  \xymatrix{
    \HH^n(X,\sh{F})
      \ar[r]\ar[dr]
    &
    \HH^n(\widehat{X},\widehat{\sh{F}})
      =
    \HH^n(\widehat{X},i^*(\sh{F}))
      \ar[d]\\
    &
    \HH^n(X,\sh{F}_k)
      =
    \HH^n(X_k,i_k^*(\sh{F})).
  }
\]
Segue imediatamente a comutatividade de \eqref{III.4.1.4.3}.
\end{env}

\begin{theorem}[4.1.5]
\label{III.4.1.5}
Seja $f:X\to Y$ um morfismo próprio de esquemas noetherianos, seja
$Y'$ um subconjunto fechado de $Y$, e ponhamos $X'=f^{-1}(Y')$.
Então, para todo $\sh{O}_X$-módulo coerente $\sh{F}$,
\[
  \RR^n\widehat{f}_*(\widehat{\sh{F}})
\]
é um $\sh{O}_{\widehat{Y}}$-módulo coerente, e os homomorfismos
$\vphi_n$, $\psi_n$ e $\rho_n$ do diagrama
\eqref{III.4.1.4.2} são isomorfismos topológicos.
\end{theorem}

\begin{proof}
Basta demonstrar que $\vphi_n$ e $\psi_n$ são isomorfismos.
Com efeito, $\RR^n f_*(\sh{F})$ é coerente \sref{III.3.2.1}, logo
seu completamento $\widehat{\RR^n f_*(\sh{F})}$ é coerente
\sref[I]{I.10.8.8}; a bicontinuidade de $\vphi_n$, $\psi_n$ e
$\rho_n$ é então automática \sref[I]{I.10.11.6}.
\end{proof}

\begin{env}[4.1.6]
\label{III.4.1.6}
\emph{Observações.}
\begin{enumerate}
  \item[(i)] Ponhamos
  \[
    \widehat{\sh{F}}_k
    =
    \widehat{\sh{F}}/
    \widehat{\sh{K}}^{\,k+1}\widehat{\sh{F}}.
  \]
  Então, imediatamente,
  \[
    \sh{F}_k=i_*(\widehat{\sh{F}}_k),
  \]
  e o homomorfismo canônico \eqref{III.4.1.3.6} é precisamente o
  homomorfismo já definido em (\hyperref[III.3.4.2.2]{3.4.2.2}):
  \[
    \RR^n\widehat{f}_*(\widehat{\sh{F}})
    \longrightarrow
    \varprojlim_k
    \RR^n\widehat{f}_*(\widehat{\sh{F}}_k).
    \tag{4.1.6.1}
    \label{III.4.1.6.1}
  \]
  Assim, o fato de que $\psi_n$ seja um isomorfismo é um caso particular
  de \sref{III.3.4.3}.
  Daremos, não obstante, mais adiante uma demonstração direta, evitando
  as delicadas considerações relativas aos limites projetivos de
  sucessões espectrais \sref[0]{0.13.7} nas quais se apoia o teorema
  geral \sref{III.3.4.3}.

  \item[(ii)] Como os $\psi_n$ são isomorfismos, dizer que os
  $\vphi_n$ são isomorfismos equivale a dizer que
  \[
    \rho_n=\psi_n^{-1}\circ\vphi_n
  \]
  são isomorfismos.
  O Teorema~\sref{III.4.1.5} afirma, em particular, que a formação de
  $\RR^n f_*$ comuta com o completamento; pode chamar-se, portanto, o
  primeiro teorema de comparação entre as teorias ``algébrica'' e
  ``formal''.
\end{enumerate}
Começaremos demonstrando a forma afim de \sref{III.4.1.5}.
\end{env}

\begin{corollary}[4.1.7]
\label{III.4.1.7}
Sob as hipóteses de \sref{III.4.1.5}, suponhamos além disso que
$Y=\Spec(A)$, onde $A$ é noetheriano, e que
$\sh{J}=\widetilde{\mathfrak{J}}$, onde $\mathfrak{J}$ é um ideal
de $A$, de modo que
\[
  \sh{F}_k=\sh{F}/\mathfrak{J}^{k+1}\sh{F}.
\]
O homomorfismo canônico
\[
  \vphi_n:
  \widehat{\HH^n(X,\sh{F})}
  \longrightarrow
  \varprojlim_k\HH^n(X,\sh{F}_k)
  \tag{4.1.7.1}
  \label{III.4.1.7.1}
\]
onde o primeiro termo é o completamento separado de
$\HH^n(X,\sh{F})$ para a topologia $\mathfrak{J}$-pré-ádica, é um
isomorfismo.
Para todo $n$, o sistema projetivo
\[
  \bigl(\HH^n(X,\sh{F}_k)\bigr)_{k\geq0}
\]
satisfaz a condição \emph{(ML)}, e o homomorfismo canônico
\[
  \psi_n:
  \HH^n(\widehat{X},\widehat{\sh{F}})
  \longrightarrow
  \varprojlim_k\HH^n(X,\sh{F}_k)
  \tag{4.1.7.2}
  \label{III.4.1.7.2}
\]

\oldpage[III]{126}

é um isomorfismo.
Por fim, a filtração sobre $\HH^n(X,\sh{F})$ definida pelos
núcleos dos homomorfismos canônicos
\[
  \HH^n(X,\sh{F})\longrightarrow\HH^n(X,\sh{F}_k)
\]
é $\mathfrak{J}$-boa \sref[0]{0.13.7.7}, e $\vphi_n$ é um
isomorfismo topológico.\footnote{A demonstração seguinte, mais simples
que a demonstração original, e o complemento relativo à filtração
de $\HH^n(X,\sh{F})$, nos foram comunicados por J.-P. Serre.}
\end{corollary}

\begin{proof}
Fixemos $n\geq0$ durante toda esta demonstração e, para simplificar,
ponhamos
\[
  \mathrm{H}=\HH^n(X,\sh{F}),
  \qquad
  \mathrm{H}_k=\HH^n(X,\sh{F}_k).
  \tag{4.1.7.3}
  \label{III.4.1.7.3}
\]
Ponhamos também
\[
  \mathrm{R}_k
  =
  \Ker(\mathrm{H}\longrightarrow\mathrm{H}_k),
  \qquad
  \text{um sub-$A$-módulo de $\mathrm{H}$.}
  \tag{4.1.7.4}
  \label{III.4.1.7.4}
\]
A sucessão exata de cohomologia
\[
  \HH^n(X,\mathfrak{J}^{k+1}\sh{F})
  \longrightarrow
  \HH^n(X,\sh{F})
  \longrightarrow
  \HH^n(X,\sh{F}_k)
  \xrightarrow{\partial}
  \HH^{n+1}(X,\mathfrak{J}^{k+1}\sh{F})
  \longrightarrow
  \HH^{n+1}(X,\sh{F})
\]
mostra que
\[
  \mathrm{R}_k
  =
  \Im\bigl(
    \HH^n(X,\mathfrak{J}^{k+1}\sh{F})
    \longrightarrow
    \HH^n(X,\sh{F})
  \bigr).
\]
Ponhamos
\[
  \begin{split}
  \mathrm{Q}_k
  &=
  \Ker\bigl(
    \HH^{n+1}(X,\mathfrak{J}^{k+1}\sh{F})
    \longrightarrow
    \HH^{n+1}(X,\sh{F})
  \bigr)
  \\
  &=
  \Im\bigl(
    \HH^n(X,\sh{F}_k)
    \longrightarrow
    \HH^{n+1}(X,\mathfrak{J}^{k+1}\sh{F})
  \bigr).
  \end{split}
  \tag{4.1.7.5}
  \label{III.4.1.7.5}
\]
Temos, portanto, a sucessão exata
\[
  0\longrightarrow\mathrm{R}_k
  \longrightarrow\mathrm{H}
  \longrightarrow\mathrm{H}_k
  \longrightarrow\mathrm{Q}_k
  \longrightarrow0.
  \tag{4.1.7.6}
  \label{III.4.1.7.6}
\]

\begin{env}[4.1.7.7]
\label{III.4.1.7.7}
Seja $x\in\mathfrak{J}^m$, onde $m\geq0$.
A multiplicação por $x$ sobre $\mathfrak{J}^k\sh{F}$ é um
homomorfismo
\[
  \mathfrak{J}^k\sh{F}
  \longrightarrow
  \mathfrak{J}^{k+m}\sh{F},
\]
e dá, portanto, um homomorfismo
\[
  \mu_{x,m}:
  \HH^n(X,\mathfrak{J}^k\sh{F})
  \longrightarrow
  \HH^n(X,\mathfrak{J}^{k+m}\sh{F}).
  \tag{4.1.7.8}
  \label{III.4.1.7.8}
\]
Seja
\[
  S=\bigoplus_{k\geq0}\mathfrak{J}^k
\]
a $A$-álgebra graduada.
Sabemos que as multiplicações $\mu_{x,m}$ dotam a
\[
  \mathrm{E}
  =
  \bigoplus_{k\geq0}
  \HH^n(X,\mathfrak{J}^k\sh{F})
\]
de uma estrutura de $S$-módulo graduado de tipo finito
\sref{III.3.3.2}; o anel graduado $S$ é noetheriano
\sref[II]{II.2.1.5}.
\end{env}

\begin{lemma}[4.1.7.9]
\label{III.4.1.7.9}
Os submódulos $\mathrm{R}_k$ de $\mathrm{H}$ definem uma
filtração $\mathfrak{J}$-boa sobre $\mathrm{H}$.
\end{lemma}

\begin{proof}
Demonstremos primeiro que
\[
  \mathfrak{J}^m\mathrm{R}_k
  \subset
  \mathrm{R}_{k+m},
  \tag{4.1.7.10}
  \label{III.4.1.7.10}
\]
onde a multiplicação sobre
$\mathrm{H}=\HH^n(X,\sh{F})$ por um elemento
$x\in\mathfrak{J}^m$ é a aplicação $\mu_{x,0}$.
Para todo $x\in\mathfrak{J}^m$, o diagrama
\[
  \xymatrix{
    \mathfrak{J}^{k+1}\sh{F}
      \ar[r]^{x}\ar[d]
    &
    \mathfrak{J}^{k+m+1}\sh{F}
      \ar[d]
    \\
    \sh{F}
      \ar[r]_{x}
    &
    \sh{F}
  }
  \tag{4.1.7.11}
  \label{III.4.1.7.11}
\]

\oldpage[III]{127}

é comutativo; suas flechas horizontais são as multiplicações por
$x$, e suas flechas verticais são as injeções canônicas.
Consequentemente, o diagrama correspondente
\[
  \xymatrix{
    \HH^n(X,\mathfrak{J}^{k+1}\sh{F})
      \ar[r]^{\mu_{x,m}}\ar[d]
    &
    \HH^n(X,\mathfrak{J}^{k+m+1}\sh{F})
      \ar[d]
    \\
    \HH^n(X,\sh{F})
      \ar[r]_{\mu_{x,0}}
    &
    \HH^n(X,\sh{F})
  }
\]
é comutativo.
A interpretação de $\mathrm{R}_k$ como a imagem de
\[
  \HH^n(X,\mathfrak{J}^{k+1}\sh{F})
  \longrightarrow
  \HH^n(X,\sh{F})
\]
demonstra, portanto, \eqref{III.4.1.7.10}.
Mostra também que o $S$-módulo graduado
\[
  \mathrm{R}=\bigoplus_{k\geq0}\mathrm{R}_k
\]
é um quociente do sub-$S$-módulo
\[
  \mathrm{M}
  =
  \bigoplus_{k\geq0}
  \HH^n(X,\mathfrak{J}^{k+1}\sh{F})
\]
de $\mathrm{E}$.
A observação anterior mostra, pois, que $\mathrm{R}$ é um
$S$-módulo de tipo finito.
Isto equivale à afirmação do lema
(Bourbaki, \emph{Alg.\ comm.}, cap.~III, \S3, n.º~1, teorema~1).
\end{proof}

\begin{env}[4.1.7.12]
\label{III.4.1.7.12}
Consideremos agora o $S$-módulo graduado
\[
  \mathrm{N}
  =
  \bigoplus_{k\geq0}
  \HH^{n+1}(X,\mathfrak{J}^{k+1}\sh{F}),
\]
definido como em \eqref{III.4.1.7.8}.
Também é um $S$-módulo de tipo finito por \sref{III.3.3.2}.
Por \eqref{III.4.1.7.5}, temos
$\mathrm{Q}_k\subset\mathrm{N}_k$ para todo $k$, e o diagrama
\eqref{III.4.1.7.11}, substituindo $n$ por $n+1$, mostra que
\[
  S_m\mathrm{Q}_k
  =
  \mathfrak{J}^m\mathrm{Q}_k
  \subset
  \mathrm{Q}_{k+m}.
\]
Em outros termos, $\mathrm{Q}$ é um sub-$S$-módulo graduado de
$\mathrm{N}$ e é, portanto, de tipo finito.
\end{env}

\begin{env}[4.1.7.13]
\label{III.4.1.7.13}
Denotemos por $\alpha_m$ a injeção canônica
\[
  \mathfrak{J}^m\longrightarrow A,
\]
que também pode escrever-se $S_m\to S_0$.
Como $\mathfrak{J}^{k+1}\sh{F}_k=0$, o $A$-módulo
$\HH^n(X,\sh{F}_k)$ é anulado por $\mathfrak{J}^{k+1}$.
Como $\mathrm{Q}_k$ é a imagem do $A$-homomorfismo
\[
  \HH^n(X,\sh{F}_k)
  \longrightarrow
  \HH^{n+1}(X,\mathfrak{J}^{k+1}\sh{F}),
\]
$\mathrm{Q}_k$ é igualmente anulado por $\mathfrak{J}^{k+1}$.
Equivalentemente, no $S$-módulo $\mathrm{Q}$,
\[
  \alpha_{k+1}(S_{k+1})\mathrm{Q}_k=0.
  \tag{4.1.7.14}
  \label{III.4.1.7.14}
\]
Como $\mathrm{Q}$ é um $S$-módulo de tipo finito, existem inteiros
$k_0$ e $h$ tais que
\[
  \mathrm{Q}_{k+h}=S_h\mathrm{Q}_k
  \qquad(k\geq k_0)
\]
\sref[II]{II.2.1.6}[ii].
Esta relação e \eqref{III.4.1.7.14} implicam que existe um inteiro
$r>0$ tal que
\[
  \alpha_r(S_r)\mathrm{Q}=0.
  \tag{4.1.7.15}
  \label{III.4.1.7.15}
\]
\end{env}

\begin{env}[4.1.7.16]
\label{III.4.1.7.16}
A injeção canônica
\[
  \mathfrak{J}^{k+m}\sh{F}
  \longrightarrow
  \mathfrak{J}^k\sh{F}
\]
induz em cohomologia um $A$-homomorfismo
\[
  \nu_m:
  \HH^{n+1}(X,\mathfrak{J}^{k+m}\sh{F})
  \longrightarrow
  \HH^{n+1}(X,\mathfrak{J}^k\sh{F}).
  \tag{4.1.7.17}
  \label{III.4.1.7.17}
\]
Para todo $x\in\mathfrak{J}^m$, temos manifestamente a
fatoração
\[
  \xymatrix@C=4.5em{
    \HH^{n+1}(X,\mathfrak{J}^k\sh{F})
      \ar[r]^{\mu_{x,m}}
    &
    \HH^{n+1}(X,\mathfrak{J}^{k+m}\sh{F})
      \ar[r]^{\nu_m}
    &
    \HH^{n+1}(X,\mathfrak{J}^k\sh{F}),
  }
  \tag{4.1.7.18}
  \label{III.4.1.7.18}
\]
cuja composta é $\mu_{x,0}$.
\end{env}

\oldpage[III]{128}

Segue-se que, para todo sub-$A$-módulo $\mathrm{P}$ de
$\HH^{n+1}(X,\mathfrak{J}^k\sh{F})$, temos no $S$-módulo
$\mathrm{N}$
\[
  \nu_m(S_m\mathrm{P})
  =
  \alpha_m(S_m)\mathrm{P}.
  \tag{4.1.7.19}
  \label{III.4.1.7.19}
\]

\begin{lemma}[4.1.7.20]
\label{III.4.1.7.20}
Existe um inteiro $m>0$ tal que
\[
  \nu_m(\mathrm{Q}_{k+m})=0
\]
para todo $k\geq k_0$.
\end{lemma}

\begin{proof}
Tomemos $m$ múltiplo de $h$ com $m\geq r$.
Como
\[
  \mathrm{Q}_{k+m}=S_m\mathrm{Q}_k
  \qquad(k\geq k_0),
\]
as relações \eqref{III.4.1.7.19} e \eqref{III.4.1.7.15} dão
\[
  \nu_m(\mathrm{Q}_{k+m})
  =
  \alpha_m(S_m)\mathrm{Q}_k
  \subset
  \alpha_r(S_r)\mathrm{Q}_k
  =
  0.
\]
\end{proof}

\begin{env}[4.1.7.21]
\label{III.4.1.7.21}
O diagrama comutativo
\[
  \xymatrix@C=2.5em{
    \HH^n(X,\sh{F})
      \ar[r]
    &
    \HH^n(X,\sh{F}_k)
      \ar[r]^-{\partial}
    &
    \HH^{n+1}(X,\mathfrak{J}^{k+1}\sh{F})
      \ar[r]
    &
    \HH^{n+1}(X,\sh{F})
    \\
    \HH^n(X,\sh{F})
      \ar[r]\ar[u]
    &
    \HH^n(X,\sh{F}_{k+m})
      \ar[r]^-{\partial}\ar[u]
    &
    \HH^{n+1}(X,\mathfrak{J}^{k+m+1}\sh{F})
      \ar[r]\ar[u]
    &
    \HH^{n+1}(X,\sh{F})
      \ar[u]
  }
\]
procede, por sua vez, do diagrama comutativo
\[
  \xymatrix{
    0\ar[r]
    &
    \mathfrak{J}^{k+1}\sh{F}\ar[r]
    &
    \sh{F}\ar[r]
    &
    \sh{F}_k\ar[r]
    &
    0
    \\
    0\ar[r]
    &
    \mathfrak{J}^{k+m+1}\sh{F}\ar[r]\ar[u]
    &
    \sh{F}\ar[r]\ar[u]
    &
    \sh{F}_{k+m}\ar[r]\ar[u]
    &
    0.
  }
\]
Aqui as flechas verticais são as aplicações canônicas.
Obtemos, portanto, o diagrama comutativo
\[
  \xymatrix@C=2.2em{
    0\ar[r]
    &
    \mathrm{R}_k\ar[r]
    &
    \mathrm{H}\ar[r]
    &
    \mathrm{H}_k\ar[r]
    &
    \mathrm{Q}_k\ar[r]
    &
    0
    \\
    0\ar[r]
    &
    \mathrm{R}_{k+m}\ar[r]\ar[u]
    &
    \mathrm{H}\ar[r]\ar@{=}[u]
    &
    \mathrm{H}_{k+m}\ar[r]\ar[u]
    &
    \mathrm{Q}_{k+m}\ar[r]\ar[u]^{\nu_m}
    &
    0,
  }
\]
de linhas exatas.
Para $k\geq k_0$, a última flecha vertical é nula por
\sref{III.4.1.7.20}.
Assim, a imagem de $\mathrm{H}_{k+m}$ em $\mathrm{H}_k$ está contida
em
\[
  \Ker(\mathrm{H}_k\longrightarrow\mathrm{Q}_k)
  =
  \Im(\mathrm{H}\longrightarrow\mathrm{H}_k).
\]
Reciprocamente, a comutatividade mostra que contém
$\Im(\mathrm{H}\to\mathrm{H}_k)$ e, portanto, que é igual a esta
imagem.
O mesmo ocorre, pois, com as imagens em $\mathrm{H}_k$ de
$\mathrm{H}_{k'}$ para $k'\geq k+m$.
Isto demonstra a condição \emph{(ML)} para o sistema projetivo
$(\mathrm{H}_k)_{k\geq0}$.

Além disso, para todo subconjunto aberto afim $U$ de $X$,
\[
  \HH^i(U,\sh{F}_k)=0
  \qquad(i>0)
\]
por \sref{III.1.3.1}, e, para $m>0$, a aplicação
\[
  \HH^0(U,\sh{F}_{k+m})
  \longrightarrow
  \HH^0(U,\sh{F}_k)
\]
é sobrejetiva \sref[I]{I.1.3.9}.
Podemos, portanto,

\oldpage[III]{129}

aplicar \sref[0]{0.13.3.1}; o homomorfismo canônico
\[
  \HH^n(\widehat{X},\widehat{\sh{F}})
  \longrightarrow
  \varprojlim_k\HH^n(X,\sh{F}_k)
\]
é bijetivo para todo $n\geq0$.

Como o sistema projetivo
$(\mathrm{H}/\mathrm{R}_k)_{k\geq0}$ é estrito, podemos passar ao
limite projetivo nas sucessões exatas
\[
  0
  \longrightarrow
  \mathrm{H}/\mathrm{R}_k
  \longrightarrow
  \mathrm{H}_k
  \longrightarrow
  \mathrm{Q}_k
  \longrightarrow
  0
  \tag{4.1.7.22}
  \label{III.4.1.7.22}
\]
por \sref[0]{0.13.2.2}.
Como $\nu_m(\mathrm{Q}_{k+m})=0$, temos
\[
  \varprojlim_k\mathrm{Q}_k=0,
\]
e, portanto, um isomorfismo topológico
\[
  \varprojlim_k(\mathrm{H}/\mathrm{R}_k)
  \isoto
  \varprojlim_k\mathrm{H}_k.
\]
Mas a filtração $(\mathrm{R}_k)$ de $\mathrm{H}$ é
$\mathfrak{J}$-boa e define, portanto, sobre $\mathrm{H}$ a
topologia $\mathfrak{J}$-pré-ádica.
Assim, $\varprojlim_k(\mathrm{H}/\mathrm{R}_k)$ é o completamento
separado de $\mathrm{H}$ para esta topologia.
Isto termina a demonstração de \sref{III.4.1.7}.
\end{env}
\end{proof}

\begin{env}[4.1.8]
\label{III.4.1.8}
Podemos terminar agora a demonstração de \sref{III.4.1.5}.
Para todo subconjunto aberto afim $V$ de $Y$,
\[
  \Gamma\bigl(
    V,\widehat{\RR^n f_*(\sh{F})}
  \bigr)
\]
é o completamento separado de
$\Gamma(V,\RR^n f_*(\sh{F}))$ para a topologia
$\mathfrak{J}$-pré-ádica, se
$\sh{J}|V=\widetilde{\mathfrak{J}}$, porque
$\RR^n f_*(\sh{F})$ é um $\sh{O}_Y$-módulo coerente
\sref[I]{I.10.8.4}.
Além disso,
\[
  \Gamma\left(
    V,\varprojlim_k\RR^n f_*(\sh{F}_k)
  \right)
  =
  \varprojlim_k
  \Gamma\bigl(V,\RR^n f_*(\sh{F}_k)\bigr)
\]
por \sref[0]{0.3.2.6}.
De \sref{III.4.1.7} e \sref{III.1.4.11} segue-se que
$\vphi_n$ é um isomorfismo topológico.
Do mesmo modo, novamente por \sref{III.1.4.11}, \sref{III.4.1.7}
mostra que o homomorfismo $\psi_{n,V}$ de \eqref{III.4.1.3.5} é um
isomorfismo.
Pela definição de
$\RR^n\widehat{f}_*(\widehat{\sh{F}})$, segue-se que $\psi_n$ é um
isomorfismo.
\end{env}

\begin{corollary}[4.1.9]
\label{III.4.1.9}
Sob as hipóteses de \sref{III.4.1.4}, para todo subconjunto aberto
afim $V$ de $Y$, o homomorfismo canônico
\[
  \HH^n\bigl(
    \widehat{X}\cap f^{-1}(V),
    \widehat{\sh{F}}
  \bigr)
  \longrightarrow
  \Gamma\bigl(
    \widehat{Y}\cap V,
    \RR^n\widehat{f}_*(\widehat{\sh{F}})
  \bigr)
\]
é bijetivo.
\end{corollary}

\begin{env}[4.1.10]
\label{III.4.1.10}
\emph{Observação.}
Seja $f:X\to Y$ um morfismo de tipo finito entre esquemas ordinários
noetherianos, e seja $\sh{F}$ um $\sh{O}_X$-módulo coerente cujo
suporte é próprio sobre $Y$ \sref[II]{II.5.4.10}.
Sabemos então \sref{III.3.2.4} que $\RR^n f_*(\sh{F})$ é um
$\sh{O}_Y$-módulo coerente para todo $n\geq0$.
Além disso, sempre podemos supor que
\[
  \sh{F}=u_*(\sh{G}),
  \qquad
  \sh{G}=u^*(\sh{F}),
\]
onde $\sh{G}$ é um $\sh{O}_Z$-módulo coerente, $Z$ é um
subesquema fechado adequado de $X$ cujo espaço subjacente é
$\Supp(\sh{F})$, e $u:Z\to X$ é a injeção canônica
\sref[I]{I.9.3.5}.
Se
\[
  \sh{G}_k
  =
  \sh{G}\otimes_{\sh{O}_Y}
  (\sh{O}_Y/\sh{J}^{k+1}),
\]
então
\[
  \sh{G}_k=u^*(\sh{F}_k),
  \qquad
  \RR^n f_*(\sh{F}_k)
  =
  \RR^n(f\circ u)_*(\sh{G}_k)
\]
e
\[
  \RR^n f_*(\sh{F})
  =
  \RR^n(f\circ u)_*(\sh{G})
\]
por \sref{III.1.3.4}; e, tendo em conta \sref[I]{I.10.9.5},
\[
  \RR^n\widehat{f}_*(\widehat{\sh{F}})
  =
  \RR^n\widehat{(f\circ u)}_*(\widehat{\sh{G}}).
\]
Podemos, portanto, aplicar \sref{III.4.1.5} a $\sh{G}$ e ao morfismo
próprio $f\circ u$.
Segue-se que, sob estas hipóteses, os resultados de
\sref{III.4.1.5} continuam válidos para $\sh{F}$ e $f$.
\end{env}


\subsection{Casos particulares e variantes}
\label{subsection:III.4.2}

A forma mais útil do teorema de comparação \sref{III.4.1.5} é a
seguinte.

\begin{proposition}[4.2.1]
\label{III.4.2.1}
Seja $Y$ um esquema localmente noetheriano, seja $f:X\to Y$ um
morfismo próprio, e seja $\sh{F}$ um $\sh{O}_X$-módulo coerente.
Então, para todo $y\in Y$ e todo $p$,
\[
  \bigl(\RR^p f_*(\sh{F})\bigr)_y
\]

\oldpage[III]{130}

é um $\sh{O}_y$-módulo de tipo finito e está, portanto, separado
para a topologia $\mathfrak{m}_y$-pré-ádica.
Além disso, existe um isomorfismo topológico canônico
\[
  \widehat{\bigl(\RR^p f_*(\sh{F})\bigr)_y}
  \isoto
  \varprojlim_k
  \HH^p\left(
    f^{-1}(y),
    \sh{F}\otimes_{\sh{O}_Y}
    (\sh{O}_y/\mathfrak{m}_y^k)
  \right).
  \tag{4.2.1.1}
  \label{III.4.2.1.1}
\]
Aqui o primeiro termo é o completamento de
$\bigl(\RR^p f_*(\sh{F})\bigr)_y$ para a topologia
$\mathfrak{m}_y$-pré-ádica.
No segundo termo, para todo $k\geq0$, $f^{-1}(y)$ se considera
como o espaço subjacente do esquema
\[
  X\times_Y\Spec(\sh{O}_y/\mathfrak{m}_y^k)
\]
\sref[I]{I.3.6.1}.
\end{proposition}

\begin{proof}
Como $\sh{O}_y$ é um anel local noetheriano e
$\bigl(\RR^p f_*(\sh{F})\bigr)_y$ é um $\sh{O}_y$-módulo de tipo
finito \sref{III.3.2.1}, a topologia $\mathfrak{m}_y$-pré-ádica sobre
este módulo é separada \sref[0]{0.7.3.5}.

Quando $Y$ é noetheriano e $y$ é fechado, as demais afirmações se
seguem de \sref{III.4.1.7}, substituindo $Y$ por uma vizinhança afim de
$y$ e tomando $Y'=\{y\}$; aqui utilizamos (G,~II,~4.9.1).

Em geral, ponhamos
\[
  Y_1=\Spec(\sh{O}_y),
  \qquad
  X_1=X\times_Y Y_1,
  \qquad
  \sh{F}_1
  =
  \sh{F}\otimes_{\sh{O}_Y}\sh{O}_{Y_1},
\]
e seja
\[
  f_1=f\times 1_{Y_1}:X_1\longrightarrow Y_1.
\]
O esquema $Y_1$ é noetheriano, o morfismo $f_1$ é próprio
\sref[II]{II.5.4.2}[iii], e $\sh{F}_1$ é coerente
\sref[0]{0.5.3.11}.
Seja $y_1$ o único ponto fechado de $Y_1$.
A proposição já demonstrada se aplica a $f_1$, $\sh{F}_1$ e
$y_1$.
Temos
\[
  \sh{O}_{y_1}=\sh{O}_y,
  \qquad
  f_1^{-1}(y_1)=f^{-1}(y)
\]
\sref[I]{I.3.6.5}.
Além disso, os esquemas
\[
  X\times_Y\Spec(\sh{O}_y/\mathfrak{m}_y^k)
  \quad\text{e}\quad
  X_1\times_{Y_1}\Spec(\sh{O}_y/\mathfrak{m}_y^k)
\]
identificam-se canonicamente \sref[I]{I.3.3.9}, assim como
\[
  \sh{F}_1
  \otimes_{\sh{O}_{Y_1}}
  (\sh{O}_y/\mathfrak{m}_y^k)
  \quad\text{e}\quad
  \sh{F}
  \otimes_{\sh{O}_Y}
  (\sh{O}_y/\mathfrak{m}_y^k)
\]
\sref[I]{I.9.1.6}.
Resta observar que
\[
  \RR^p f_{1*}(\sh{F}_1)
  \simeq
  \RR^p f_*(\sh{F})
  \otimes_{\sh{O}_Y}
  \sh{O}_{Y_1}.
\]
Isto se segue de \sref{III.1.4.15}, pois o morfismo local
$\Spec(\sh{O}_y)\to Y$ é plano
\sref[0]{0.6.7.1}, \sref[I]{I.2.4.2}.
\end{proof}

O corolário seguinte utiliza a terminologia da teoria da
dimensão (capítulo~IV) e não se aplicará antes desse capítulo.

\begin{corollary}[4.2.2]
\label{III.4.2.2}
Seja $Y$ um esquema localmente noetheriano, seja $f:X\to Y$ um
morfismo próprio, seja $y\in Y$, e seja $r$ a dimensão de
$f^{-1}(y)$.
Então, para todo $\sh{O}_X$-módulo coerente $\sh{F}$, os feixes
$\RR^p f_*(\sh{F})$ se anulam em uma vizinhança de $y$ para todo
$p>r$.
\end{corollary}

\begin{proof}
Para todo $k$, temos
\[
  \HH^p\left(
    f^{-1}(y),
    \sh{F}\otimes_{\sh{O}_Y}
    (\sh{O}_y/\mathfrak{m}_y^k)
  \right)
  =
  0
\]
por (G,~II,~4.15.2).
De \sref{III.4.2.1} segue-se que o completamento separado de
$\bigl(\RR^p f_*(\sh{F})\bigr)_y$ para a topologia
$\mathfrak{m}_y$-pré-ádica é nulo.
Como esta topologia está separada, temos também
\[
  \bigl(\RR^p f_*(\sh{F})\bigr)_y=0.
\]
A conclusão se segue porque $\RR^p f_*(\sh{F})$ é coerente
\sref[0]{0.5.2.2}.
\end{proof}

\begin{env}[4.2.3]
\label{III.4.2.3}
O resultado \sref{III.4.2.1} é utilizado especialmente quando $p=0$.
Obtemos o corolário seguinte.
\end{env}

\begin{corollary}[4.2.4]
\label{III.4.2.4}
Sob as hipóteses de \sref{III.4.2.1}, existe um isomorfismo
topológico canônico
\[
  \widehat{\bigl(f_*(\sh{F})\bigr)_y}
  \isoto
  \varprojlim_k
  \Gamma\left(
    f^{-1}(y),
    \sh{F}_y/\mathfrak{m}_y^k\sh{F}_y
  \right).
\]
\end{corollary}


\subsection{Teorema de conexidade de Zariski}
\label{subsection:III.4.3}
\label{III.4.3}

Os resultados desta seção e da seguinte generalizam teoremas
bem conhecidos de Zariski, e todos podem ser deduzidos de
\sref{III.4.2.4}. São consequências do teorema seguinte.

\begin{theorem}[4.3.1]
\label{III.4.3.1}
\textnormal{(Teorema de conexidade).}
Seja $Y$ um esquema localmente noetheriano e seja $f:X\to Y$ um
morfismo próprio. Então
\[
  \sh{A}(X)=f_*(\sh{O}_X)
\]
é uma $\sh{O}_Y$-álgebra coerente.

\oldpage[III]{131}

Seja $Y'$ o $Y$-esquema finito tal que
$\sh{A}(Y')=\sh{A}(X)$; está determinado salvo $Y$-isomorfismo
\sref[II]{II.1.3.1}, \sref[II]{II.6.1.3}.
Se $f'=\sh{A}(e)$ é o $Y$-morfismo $X\to Y'$ deduzido do
isomorfismo identidade
\[
  e:\sh{A}(Y')\isoto\sh{A}(X)
\]
\sref[II]{II.1.2.7}, então $f'$ é próprio,
$f'_*(\sh{O}_X)$ é isomorfo a $\sh{O}_{Y'}$, e toda fibra
$f'^{-1}(y')$, $y'\in Y'$, é conexa e não vazia.
\end{theorem}

\begin{proof}
Seja $g:Y'\to Y$ o morfismo estrutural. Para demonstrar que o
homomorfismo
\[
  \theta:\sh{O}_{Y'}\longrightarrow f'_*(\sh{O}_X)
\]
que intervém na definição de $f'$ é bijetivo, basta, como
$Y'$ é afim sobre $Y$, demonstrar que
\[
  g_*(\theta):
  g_*(\sh{O}_{Y'})
  \longrightarrow
  g_*\bigl(f'_*(\sh{O}_X)\bigr)
  =
  f_*(\sh{O}_X)
\]
é a identidade \sref[II]{II.1.4.2}. Isto se segue das
definições, pois
$g_*(\sh{O}_{Y'})=\sh{A}(Y')$ e
$f_*(\sh{O}_X)=\sh{A}(X)$.
A coerência de $\sh{A}(X)$ é um caso particular do teorema de
finitude \sref{III.3.2.1}. Como $f$ é próprio e $g$ é separado,
$f'$ é próprio \sref[II]{II.5.4.3}[i]. Resta demonstrar
\sref{III.4.3.2}.
\end{proof}

\begin{corollary}[4.3.2]
\label{III.4.3.2}
Sob as hipóteses de \sref{III.4.3.1}, suponhamos além disso que
$f_*(\sh{O}_X)$ é isomorfo a $\sh{O}_Y$. Então toda fibra
$f^{-1}(y)$ de $f$ é conexa e não vazia.
\end{corollary}

\begin{proof}
A hipótese de que $f_*(\sh{O}_X)$ é isomorfo a $\sh{O}_Y$ implica
que $f$ já é dominante, logo sobrejetivo porque $f$ é uma
aplicação fechada. Como em \sref{III.4.2.1}, podemos reduzir-nos ao
caso em que $y$ é fechado em $Y$. Como $f^{-1}(y)$ é um espaço
noetheriano, tem um número finito de componentes conexas; estas são
também as componentes conexas do espaço subjacente do
completamento $\widehat{X}$ ao longo de $f^{-1}(y)$.

Se $Z_i$, $1\leq i\leq n$, são estas componentes conexas, então
\[
  \Gamma(\widehat{X},\sh{O}_{\widehat{X}})
  =
  \prod_{i=1}^{n}\Gamma(Z_i,\sh{O}_{\widehat{X}}),
\]
e nenhum dos anéis do membro direito é nulo, pois a seção
unidade não é nula em nenhum ponto de $\widehat{X}$. Aplicando
\sref{III.4.1.5} a $\sh{F}=\sh{O}_X$, cujo completamento ao longo de
$f^{-1}(y)$ é $\sh{O}_{\widehat{X}}$, vê-se que
$\Gamma(\widehat{X},\sh{O}_{\widehat{X}})$ é isomorfo ao
completamento $\mathfrak{m}_y$-ádico separado
$\widehat{\sh{O}}_y$ do anel local $\sh{O}_y$. É, portanto, um
anel local e não pode ser um produto direto de vários anéis não
nulos. Assim, $n=1$.
\end{proof}

\begin{corollary}[4.3.3]
\label{III.4.3.3}
Sob as hipóteses de \sref{III.4.3.1}, para todo $y\in Y$ o conjunto
de componentes conexas da fibra $f^{-1}(y)$ está em correspondência
bijetiva com o conjunto finito de pontos da fibra
$g^{-1}(y)$, onde $g:Y'\to Y$ é o morfismo estrutural; em outros
termos, corresponde ao conjunto de ideais maximais de
$\bigl(f_*(\sh{O}_X)\bigr)_y$.
\end{corollary}

\begin{proof}
Como $Y'$ é finito sobre $Y$, o espaço $g^{-1}(y)$ é finito e
discreto \sref[II]{II.6.1.7}. Como
\[
  f^{-1}(y)=f'^{-1}\bigl(g^{-1}(y)\bigr),
\]
a afirmação se segue desta observação e de
\sref{III.4.3.1}.
\end{proof}

Isto dá uma interpretação notável do $Y$-esquema $Y'$ definido em
\sref{III.4.3.1}. A fatoração
$f=g\circ f'$ do morfismo próprio $f$ é análoga à fatoração
obtida por K.~Stein para aplicações holomorfas de espaços
analíticos, e será chamada daqui em diante de \emph{fatoração de Stein}
de $f$.

\begin{remark}[4.3.4]
\label{III.4.3.4}
Seja $k$ uma extensão do corpo $\kres(y)$. Se
\[
  f^{-1}(y)\otimes_{\kres(y)}k
  =
  X\times_Y\Spec(k)
\]
é conexo, então $f^{-1}(y)$ também é conexo, pois é a imagem
do primeiro por um morfismo de projeção
\sref[I]{I.3.4.7}. Para um morfismo $f:X\to Y$ de esquemas e um ponto
$y\in Y$, diremos que a fibra $f^{-1}(y)$ é
\emph{geometricamente conexa} se, para toda extensão $k$ de
$\kres(y)$, o esquema
$f^{-1}(y)\otimes_{\kres(y)}k=X\times_Y\Spec(k)$ é conexo.

\oldpage[III]{132}

Sob as hipóteses de \sref{III.4.3.2}, a conclusão pode
ser reforçada: as fibras $f^{-1}(y)$ são geometricamente conexas.
Com efeito, para toda extensão $k$ de $\kres(y)$ existem um anel
local noetheriano $A$ e um homomorfismo local
\[
  \varphi:\sh{O}_y\longrightarrow A
\]
tal que $A$ é um $\sh{O}_y$-módulo plano e o corpo residual de
$A$ é $\kres(y)$-isomorfo a $k$ \sref[0]{0.10.3.1}.
Seja $Y_1=\Spec(A)$, e seja $h:Y_1\to Y$ o morfismo local
correspondente a $\varphi$, que leva o único ponto fechado $y_1$ de
$Y_1$ a $y$ \sref[I]{I.2.4.1}. Ponhamos
\[
  X_1=X\times_Y Y_1,
  \qquad
  f_1=f\times 1_{Y_1}.
\]
O esquema $Y_1$ é noetheriano, $f_1$ é próprio
\sref[II]{II.5.4.2}, e $f_1^{-1}(y_1)$ é um
$\kres(y_1)$-esquema isomorfo a $X\times_Y\Spec(k)$. Basta, portanto,
demonstrar que
\[
  f_{1*}(\sh{O}_{X_1})=\sh{O}_{Y_1},
\]
de modo que \sref{III.4.3.2} possa ser aplicado a $f_1$. Agora $h$ é
plano por \sref[I]{I.2.4.2}; portanto, por \sref{III.1.4.15} com
$q=0$,
\[
  f_{1*}(\sh{O}_{X_1})
  =
  h^*\bigl(f_*(\sh{O}_X)\bigr)
  =
  h^*(\sh{O}_Y)
  =
  \sh{O}_{Y_1}.
\]

No caso geral de \sref{III.4.3.1}, o mesmo argumento mostra,
com a notação desse teorema, que
\[
  f_{1*}(\sh{O}_{X_1})
  =
  h^*\bigl(g_*(\sh{O}_{Y'})\bigr),
\]
e que a fatoração de Stein $f_1=g_1\circ f'_1$ de $f_1$
satisfaz
\[
  g_1=g\times 1_{Y_1}
\]
\sref[II]{II.1.5.2}, sendo o $Y_1$-esquema finito correspondente
$Y'_1=Y'\times_Y Y_1$. Tendo em conta a transitividade das
fibras \sref[I]{I.3.6.4}, segue-se de
\sref{III.4.3.3} que o número de componentes conexas de
$f_1^{-1}(y_1)$ é o número de elementos de
\[
  g_1^{-1}(y_1)
  =
  g^{-1}(y)\otimes_{\kres(y)}k.
\]
Se $k$ for escolhido algebricamente fechado, este número será independente
da extensão algebricamente fechada escolhida e igual ao
\emph{número geométrico de pontos} de $g^{-1}(y)$
\sref[I]{I.6.4.7}, ou, equivalentemente, a
\[
  \sum_i[\kres(y'_i):\kres(y)]_s,
\]
onde os $y'_i$ percorrem o conjunto finito $g^{-1}(y)$. Este número
chama-se \emph{número geométrico de componentes conexas} de
$f^{-1}(y)$. Os corpos $\kres(y'_i)$ são precisamente os corpos
residuais do anel semilocal $\bigl(f_*(\sh{O}_X)\bigr)_y$.
\end{remark}

\begin{proposition}[4.3.5]
\label{III.4.3.5}
Sejam $X$ e $Y$ esquemas íntegros localmente noetherianos, e seja
$f:X\to Y$ um morfismo próprio dominante. Para todo $y\in Y$, o
número de componentes conexas de $f^{-1}(y)$ é no máximo o número
de ideais maximais do fecho integral $\sh{O}'_y$ de
$\sh{O}_y$ no corpo de funções racionais $\mathrm{R}(X)$.
\end{proposition}

\begin{proof}
Para todo aberto $U$ de $Y$,
\[
  \Gamma\bigl(U,f_*(\sh{O}_X)\bigr)
  =
  \Gamma\bigl(f^{-1}(U),\sh{O}_X\bigr)
\]
é a interseção dos anéis locais $\sh{O}_x$ para
$x\in f^{-1}(U)$ \sref[I]{I.8.2.1.1}. Segue-se imediatamente que
$\bigl(f_*(\sh{O}_X)\bigr)_y$ é um subanel de $\mathrm{R}(X)$ que
contém $\sh{O}_y$. Como $f_*(\sh{O}_X)$ é um
$\sh{O}_Y$-módulo coerente, $\bigl(f_*(\sh{O}_X)\bigr)_y$ é um
$\sh{O}_y$-módulo finito e, portanto, está contido em $\sh{O}'_y$.
Todo ideal maximal de tal anel é a interseção deste anel com
um ideal maximal de $\sh{O}'_y$
\cite[vol.~I, pp.~257 and 259]{I-13}, o que demonstra a proposição.
\end{proof}

\begin{definition}[4.3.6]
\label{III.4.3.6}
Um anel local íntegro se diz \emph{unirrama} se seu fecho integral
é novamente um anel local. Um ponto $y$ de um esquema íntegro
$Y$ se diz \emph{unirrama} se o anel local $\sh{O}_y$ é
unirrama; este é, em particular, o caso quando $Y$ é normal em
$y$.
\end{definition}

Seja $A$ um anel local íntegro e seja $K$ seu corpo de frações. O
anel $A$ é unirrama se e somente se todo subanel $A_1$ de $K$ que
contenha $A$ e seja uma $A$-álgebra finita é local. Com efeito, seja
$A'$ o fecho integral de $A$. Pelo primeiro teorema de
Cohen--Seidenberg (Bourbaki, \emph{Alg\`ebre commutative}, capítulo~V,
\S2, n.º~1, teorema~1), todo ideal maximal de $A_1$ é a contração
a $A_1$ de um ideal maximal de $A'$; assim, se $A'$ é local, também
o é $A_1$.

\oldpage[III]{133}

Reciprocamente, $A'$ é o limite indutivo da família dirigida de
$A$-subálgebras finitas $A_\alpha$ de $A'$. Se cada $A_\alpha$ é
local, então, para $A_\alpha\subset A_\beta$, o ideal maximal de
$A_\alpha$ é a contração do ideal maximal de $A_\beta$, pelo
mesmo argumento; portanto, $A'$ é local
\sref[0]{0.10.3.1.3}.

Observemos também que, se o completamento de um anel local
noetheriano $A$ é íntegro ---propriedade que se exprime dizendo que
$A$ é \emph{analiticamente íntegro}---, então $A$ é unirrama.
Seja $\mathfrak{m}$ o ideal maximal de $A$, seja $K$ seu corpo de
frações, e seja $K'$ o corpo de frações de $\widehat{A}$;
então
\[
  K'=K\otimes_A\widehat{A}.
\]
Seja $A'_F$ uma $A$-subálgebra finita de $K$. O subanel $B_F$ de
$K'$ gerado por $\widehat{A}$ e $A'_F$ é isomorfo a
\[
  A'_F\otimes_A\widehat{A};
\]
é um $\widehat{A}$-módulo finito e é o completamento de $A'_F$
para a topologia $\mathfrak{m}$-ádica
\sref[0]{0.7.3.3}, \sref[0]{0.7.3.6}. Como $A'_F$ é semilocal
(Bourbaki, \emph{Alg\`ebre commutative}, capítulo~IV, \S2, n.º~5,
corolário~3 da proposição~9), enquanto seu completamento é
íntegro, $A'_F$ só pode ter um ideal maximal $\mathfrak{m}'_F$, e
$\mathfrak{m}'_F\cap A=\mathfrak{m}$. Isto demonstra a afirmação.

\begin{corollary}[4.3.7]
\label{III.4.3.7}
Sob as hipóteses de \sref{III.4.3.5}, suponhamos que o fecho
algébrica de $\mathrm{R}(Y)$ em $\mathrm{R}(X)$ tem grau separável
$n$ e que $y\in Y$ é unirrama. Então a fibra $f^{-1}(y)$ tem no
máximo $n$ componentes conexas. Em particular, se o fecho
algébrica de $\mathrm{R}(Y)$ em $\mathrm{R}(X)$ é radicial sobre
$\mathrm{R}(Y)$, então $f^{-1}(y)$ é conexa.
\end{corollary}

\begin{proof}
Seja $\sh{O}'_y$ o fecho integral de $\sh{O}_y$, e seja
$\sh{O}''_y$ o fecho integral de $\sh{O}_y$ em $\mathrm{R}(X)$.
Esta última é também o fecho integral de $\sh{O}'_y$ em
$\mathrm{R}(X)$. Se $\sh{O}'_y$ é local, então $\sh{O}''_y$ é
um anel semilocal que tem no máximo $n$ ideais maximais
\cite[vol.~I, p.~289, Theorem~22]{I-13}.

Este corolário é essencialmente a forma na qual Zariski enuncia seu
``teorema de conexidade'' para esquemas algébricos.
\end{proof}

\begin{remark}[4.3.8]
\label{III.4.3.8}
Se, além das hipóteses de \sref{III.4.3.7}, $Y$ é normal em
$y$, então a fibra $f^{-1}(y)$ é geometricamente conexa: com a
notação de \sref{III.4.3.4}, $g^{-1}(y)$ é um único ponto $y'$, e
$\kres(y')$ é radicial sobre $\kres(y)$.
\end{remark}

\begin{definition}[4.3.9]
\label{III.4.3.9}
Seja $Y$ um esquema localmente noetheriano. Um morfismo de tipo finito
$f:X\to Y$ se diz \emph{universalmente aberto} se, para todo esquema
irredutível localmente noetheriano $Y'$ e todo morfismo dominante
$g:Y'\to Y$, toda componente irredutível de
$X'=X\times_Y Y'$ domina $Y'$.
\end{definition}

Se $Y$ é irredutível, isto significa que, se $\eta$ e $\eta'$ são
os pontos genéricos de $Y$ e $Y'$, respectivamente, e se
$f'=f\times 1_{Y'}$, toda componente irredutível de $X'$ tem seu
ponto genérico em $f'^{-1}(\eta')$. Consequentemente, para todo
aberto $U$ de $Y$, o morfismo restrito $f^{-1}(U)\to U$ é
universalmente aberto.

\begin{corollary}[4.3.10]
\label{III.4.3.10}
Sejam $X$ e $Y$ esquemas íntegros localmente noetherianos, e seja
$f:X\to Y$ próprio, dominante e universalmente aberto. Se o fecho
algébrica de $\mathrm{R}(Y)$ em $\mathrm{R}(X)$ é radicial sobre
$\mathrm{R}(Y)$, então toda fibra $f^{-1}(y)$, $y\in Y$, é
geometricamente conexa.
\end{corollary}

\begin{proof}
Podemos reduzir-nos ao caso $Y=\Spec(B)$, onde $B$ é um anel
íntegro noetheriano. Por \sref[II]{II.7.1.7}, existe um anel local
noetheriano integralmente fechado $A$ que domina $\sh{O}_y$ e tem
$\mathrm{R}(Y)$ como corpo de frações. Ponhamos $Y'=\Spec(A)$, e
seja $h:Y'\to Y$ o morfismo correspondente à injeção canônica
$B\to A$. É birracional, logo dominante; além disso, se $y'$ é o
único ponto fechado de $Y'$, então $h(y')=y$.

Ponhamos
\[
  X'=X\times_Y Y',
  \qquad
  f'=f\times 1_{Y'}.
\]
Sejam $\eta$, $\eta'$ e $\xi$ os pontos genéricos de $Y$, $Y'$ e
$X$, respectivamente. Então $f(\xi)=\eta$ e
$h^{-1}(\eta)=\{\eta'\}$; além disso,
$\kres(\eta)=\kres(\eta')=\mathrm{R}(Y)$. Assim, $f'^{-1}(\eta')$ é
isomorfo a $f^{-1}(\eta)$ \sref[I]{I.3.6.4}, e tem um único ponto
genérico, pois $\xi$ é o ponto genérico de
$f^{-1}(\eta)$ \sref[0]{0.2.1.8}. Pela abertura universal, toda
componente irredutível de $X'$ tem seu ponto genérico em
$f'^{-1}(\eta')$. Portanto, $X'$ é irredutível; seu ponto genérico
$\xi'$ é o de $f'^{-1}(\eta')$, e
$\kres(\xi')=\kres(\xi)$.

\oldpage[III]{134}

Ponhamos $X''=X'_{\mathrm{red}}$ e $f''=f'_{\mathrm{red}}$. Então
$X''$ é íntegro e noetheriano, $f''$ é próprio
\sref[II]{II.5.4.6}, e os espaços subjacentes das fibras
$f'^{-1}(y')$ e $f''^{-1}(y')$ são iguais. Além disso,
\[
  \mathrm{R}(X'')=\kres(\xi')=\mathrm{R}(X).
\]
Assim, $f''$ satisfaz as hipóteses de \sref{III.4.3.8}, e
$f''^{-1}(y')$ é geometricamente conexa.

Seja agora $k$ uma extensão qualquer de $\kres(y)$. Existe uma
extensão $k_1$ de $\kres(y)$ na qual tanto $\kres(y')$ como $k$
podem ser vistos como subextensões (Bourbaki, \emph{Alg\`ebre},
capítulo~V, \S4, proposição~2). Pelo parágrafo anterior,
\[
  f''^{-1}(y')\times_{Y'}\Spec(k_1)
\]
é conexo. Tem o mesmo esquema reduzido associado que
$f'^{-1}(y')\times_{Y'}\Spec(k_1)$
\sref[I]{I.5.1.8}, logo este último é conexo; é isomorfo a
$f^{-1}(y)\times_Y\Spec(k_1)$ \sref[I]{I.3.6.4}. Portanto,
$f^{-1}(y)\times_Y\Spec(k)$ é conexo, e a observação do início
de \sref{III.4.3.4} termina a demonstração.
\end{proof}

\begin{remark}[4.3.11]
\label{III.4.3.11}
\emph{(i)}
O argumento anterior se deve, no essencial, a Zariski
\cite{I-20}, salvo que, em seu contexto, pode-se tomar como $A$ o
fecho integral de $\sh{O}_y$, pois este é um anel noetheriano para
os anéis locais da geometria algébrica clássica. Zariski
demonstra também que, se $Y$ é a variedade de Chow de um espaço
projetivo $\mathbf{P}_k^r$ sobre um corpo $k$, e $X$ é o
subconjunto fechado de $\mathbf{P}_k^r\times_kY$ que define a
correspondência de Chow entre $\mathbf{P}_k^r$ e $Y$, então a
projeção $X\to Y$ é universalmente aberta (\emph{loc.\ cit.},
lema da p.~82).
Parece que esta é a única propriedade formal das ``coordenadas de
Chow'' utilizada em certas aplicações. Em tal situação é, por
isso, útil substituir a linguagem da especialização de ciclos no
espaço projetivo pela linguagem das fibras de um morfismo próprio,
eventualmente suposto universalmente aberto ou submetido a condições
análogas de regularidade local.

\emph{(ii)}
No capítulo~IV veremos que um morfismo universalmente aberto
$f:X\to Y$ pode definir-se equivalentemente exigindo que, para todo
morfismo $Z\to Y$, o morfismo obtido por mudança de base
$f_{(Z)}:X_{(Z)}\to Z$ seja aberto. Também pode demonstrar-se que, se
$f$ satisfaz as hipóteses de \sref{III.4.3.10}, e se $y$ é uma
especialização de $y'$, então o número geométrico de componentes
conexas de $f^{-1}(y)$ é no máximo o de $f^{-1}(y')$.
\end{remark}

\begin{corollary}[4.3.12]
\label{III.4.3.12}
Sob as hipóteses de \sref{III.4.3.5}, suponhamos além disso que
$\mathrm{R}(Y)$ é algebricamente fechado em $\mathrm{R}(X)$, e seja
$y$ um ponto normal de $Y$. Então $f^{-1}(y)$ é geometricamente
conexa, e existe uma vizinhança aberta $U$ de $y$ em $Y$ tal que
\[
  f_*\bigl(\sh{O}_X|_{f^{-1}(U)}\bigr)
  \simeq
  \sh{O}_Y|_U.
\]
Em particular, se $Y$ é normal e $\mathrm{R}(Y)$ é algebricamente
fechado em $\mathrm{R}(X)$, então
$f_*(\sh{O}_X)\simeq\sh{O}_Y$.
\end{corollary}

\begin{proof}
A afirmação relativa a $f^{-1}(y)$ é um caso particular de
\sref{III.4.3.8}.

\oldpage[III]{135}

Se
\[
  X\xrightarrow{f'}Y'\xrightarrow{g}Y
\]
é a fatoração de Stein de $f$ \sref{III.4.3.3}, segue-se que
$g^{-1}(y)$ consiste em um único ponto $y'$. Além disso,
\[
  \sh{O}_y
  \subset
  \sh{O}_{y'}
  =
  \bigl(f_*(\sh{O}_X)\bigr)_y
  \subset
  \mathrm{R}(X).
\]
Como $\sh{O}_{y'}$ é finito sobre $\sh{O}_y$, e portanto
\emph{a fortiori} algébrico sobre $\mathrm{R}(Y)$ dentro do corpo
de frações pertinente, a hipótese obriga que esteja contido em
$\mathrm{R}(Y)$. Como $y$ é normal, necessariamente
$\sh{O}_{y'}=\sh{O}_y$. Assim, $g$ é um isomorfismo local em $y'$
\sref[I]{I.6.5.4}, o que demonstra a primeira parte. A segunda
afirmação se segue da primeira: a hipótese adicional faz com que $g$
seja bijetivo e um isomorfismo local em uma vizinhança de cada ponto de
$Y'$, logo um isomorfismo.
\end{proof}

O fato de que \sref{III.4.3.7} esteja disponível para esquemas fornece,
por exemplo, a aplicação seguinte.

\begin{proposition}[4.3.13]
\label{III.4.3.13}
Seja $A$ um anel local noetheriano unirrama, seja $\mathfrak{a}$ um
ideal de definição de $A$, e sejam
\[
  A_0=A/\mathfrak{a},
  \qquad
  S=\operatorname{gr}_{\mathfrak{a}}(A)
\]
o anel graduado associado à filtração $\mathfrak{a}$-pré-ádica.
Suponhamos que $S$ é uma $A_0$-álgebra graduada gerada por $S_1$,
onde $S_1$ é um $A_0$-módulo finito. Então $\Proj(S)$ é um
$A_0$-esquema conexo.
\end{proposition}

\begin{proof}
Seja $\mathfrak{m}$ o ideal maximal de $A$. O esquema
$Y=\Spec(A)$ é íntegro, e o ponto $y$ correspondente a
$\mathfrak{m}$ é seu único ponto fechado. Para certo inteiro $p$,
\[
  \mathfrak{m}^p\subset\mathfrak{a}\subset\mathfrak{m},
\]
de modo que $V(\mathfrak{a})=\{\mathfrak{m}\}$. Ponhamos
\[
  S'=\bigoplus_{n\geq0}\mathfrak{a}^n
  \qquad\text{e}\qquad
  X=\Proj(S').
\]
Este é o $Y$-esquema obtido pela explosão do ideal
$\mathfrak{a}$. É íntegro, e o morfismo estrutural
$f:X\to Y$ é birracional \sref[II]{II.8.1.4} e projetivo.
Portanto, \sref{III.4.3.7} se aplica e mostra que $f^{-1}(y)$ é
conexo. Mas o espaço subjacente de $f^{-1}(y)$ é o de
\[
  \Proj(S'\otimes_A A_0)
\]
\sref[I]{I.3.6.1}, \sref[II]{II.2.8.10}; como
$S'\otimes_AA_0=S$ por definição, segue-se a proposição.
\end{proof}


\subsection{O ``teorema principal'' de Zariski}
\label{subsection:III.4.4}

\begin{proposition}[4.4.1]
\label{III.4.4.1}
Seja $Y$ um esquema localmente noetheriano e seja $f:X\to Y$ um
morfismo próprio. Seja $X'$ o conjunto de pontos $x\in X$ que estão
isolados em sua fibra $f^{-1}(f(x))$. Então $X'$ é aberto em $X$.
Se
\[
  f=g\circ f'
\]
é a fatoração de Stein de $f$ \sref{III.4.3.3}, a restrição
de $f'$ a $X'$ é um isomorfismo de $X'$ sobre o subesquema induzido
em um subconjunto aberto $U$ de $Y'$, e
$X'=f'^{-1}(U)$.
\end{proposition}

\begin{proof}
Como $g^{-1}(f(x))$ é finito e discreto
\sref{III.4.3.3}, \sref[II]{II.6.1.7}, o ponto $x$ está isolado em
$f^{-1}(f(x))$ se e somente se está isolado em
$f'^{-1}(f'(x))$. Podemos, portanto, reduzir-nos ao caso $f'=f$, de
modo que $f_*(\sh{O}_X)=\sh{O}_Y$. Se $x\in X'$, a fibra
$f^{-1}(f(x))$, que é conexa por \sref{III.4.3.2}, deve constar
então do único ponto $x$.

Como $f$ é fechado, para toda vizinhança aberta $V$ de $x$ em $X$, o
conjunto fechado $f(X-V)$ não contém $y=f(x)$. Se $U$ é seu
complementar em $Y$, então $f^{-1}(U)\subset V$. Assim, as
imagens inversas por $f$ de um sistema fundamental de vizinhanças
abertas de $y$ formam um sistema fundamental de vizinhanças abertas de
$x$. A igualdade $f_*(\sh{O}_X)=\sh{O}_Y$ e a definição da
imagem direta de um feixe \sref[0]{0.3.4.1}, \sref{III.4.2.1} implicam
agora que, se $f=(\psi,\theta)$, o homomorfismo
\[
  \theta_y^\#:\sh{O}_y\longrightarrow\sh{O}_x
\]
é um isomorfismo. Portanto, existem vizinhanças abertas $V$ de $x$ e
$U$ de $y$ tais que a restrição de $f$ é um isomorfismo
$V\isoto U$ \sref[I]{I.6.5.4}.

\oldpage[III]{136}

Pelo parágrafo anterior podemos dispor além disso que
$f^{-1}(U)=V$. Segue-se imediatamente que $V\subset X'$, o que
termina a demonstração.
\end{proof}

A proposição seguinte foi demonstrada por Chevalley para esquemas
algébricos.

\begin{proposition}[4.4.2]
\label{III.4.4.2}
Seja $Y$ um esquema localmente noetheriano e seja $f:X\to Y$ um
morfismo. As condições seguintes são equivalentes:
\begin{enumerate}
  \item[{\rm a)}] $f$ é finito;
  \item[{\rm b)}] $f$ é afim e próprio;
  \item[{\rm c)}] $f$ é próprio e, para todo $y\in Y$, o conjunto
    $f^{-1}(y)$ é finito.
\end{enumerate}
\end{proposition}

\begin{proof}
A condição~a) implica b) por \sref[II]{II.6.1.2} e
\sref[II]{II.6.1.11}. Se $f$ é próprio e afim, também o é
\[
  f^{-1}(y)\longrightarrow\Spec(\kres(y))
\]
\sref[II]{II.1.6.2}, \sref[II]{II.5.4.2}. O teorema de finitude
\sref{III.3.2.1}, aplicado ao feixe estrutural de $f^{-1}(y)$, mostra
que $f^{-1}(y)=\Spec(A)$ para uma $\kres(y)$-álgebra finita $A$. Assim,
$f^{-1}(y)$ é um conjunto finito \sref[II]{II.6.1.7}, e b) implica
c).

Por fim, se c) for satisfeita, a fibra $f^{-1}(y)$ é um esquema
algébrico sobre $\kres(y)$ cujo conjunto subjacente é finito e é,
portanto, discreto \sref[I]{I.6.4.4}. Com a notação de
\sref{III.4.4.1}, temos $X'=X$, logo
$f':X\to Y'$ é um isomorfismo. Como $g$ é finito, c) implica a).
\end{proof}

\begin{theorem}[4.4.3]
\label{III.4.4.3}
\textnormal{(``Teorema principal'' de Zariski).}
Seja $Y$ um esquema noetheriano, seja $f:X\to Y$ um morfismo
quase-projetivo, e seja $X'$ o conjunto de pontos $x\in X$ que estão
isolados em sua fibra $f^{-1}(f(x))$. Então $X'$ é aberto em $X$,
e o subesquema induzido por $X$ sobre $X'$ é isomorfo a um
subesquema induzido sobre um subconjunto aberto de um $Y$-esquema
$Y'$ finito sobre $Y$.
\end{theorem}

\begin{proof}
Existe um $Y$-esquema projetivo $Z$ tal que $X$ é
$Y$-isomorfo ao subesquema induzido sobre um subconjunto aberto de
$Z$ \sref[II]{II.5.3.2}, \sref[II]{II.5.5.1}. Reduzimo-nos assim ao
caso em que $f$ é projetivo, logo próprio \sref[II]{II.5.5.3}, e o
teorema se segue imediatamente de
\sref{III.4.4.1}.
\end{proof}

\begin{remark}[4.4.4]
\label{III.4.4.4}
Se $X$ é reduzido (respectivamente, se $X$ é irredutível e $X'$ não
é vazio), pode exigir-se em \sref{III.4.4.3} que $Y'$ seja reduzido
(respectivamente, irredutível). Com efeito, $Y'$ pode ser substituído pelo
fecho esquemático $\overline{X'}$ de $X'$ em $Y'$
\sref[I]{I.9.5.11}, \sref[II]{II.6.1.5}. Se $X'$ é reduzido, também
o é $\overline{X'}$ \sref[I]{I.9.5.9}; se $X'$ não é vazio e é
irredutível, então $\overline{X'}$ também é irredutível.
\end{remark}

\begin{corollary}[4.4.5]
\label{III.4.4.5}
Seja $Y$ um esquema localmente noetheriano, seja $f:X\to Y$ um
morfismo de tipo finito, e seja $x$ um ponto de $X$ isolado em sua fibra
$f^{-1}(f(x))$. Então existe uma vizinhança aberta de $x$ em $X$
isomorfa a um subconjunto aberto de um $Y$-esquema finito sobre $Y$.
\end{corollary}

\begin{proof}
Ponhamos $y=f(x)$, seja $U$ uma vizinhança aberta afim de $y$ em $Y$, e
seja $V$ uma vizinhança aberta afim de $x$ em $X$, contida em
$f^{-1}(U)$. Como $Y$ é separado, a inclusão $U\to Y$ é afim
\sref[II]{II.1.6.3}; e como $V$ é afim sobre $U$ (\emph{ibid.}), a
restrição de $f$ a $V$ é um morfismo afim $V\to Y$
\sref[II]{II.1.6.2}. A fortiori, esta restrição é
quase-projetiva, pois é de tipo finito
\sref[I]{I.6.3.5}, \sref[II]{II.5.3.4}. Resta aplicar
\sref{III.4.4.3} a esta restrição.
\end{proof}

O Corolário \sref{III.4.4.5} pode ser enunciado na linguagem da
álgebra comutativa como segue.

\begin{corollary}[4.4.6]
\label{III.4.4.6}
Seja $A$ um anel noetheriano, seja $B$ uma $A$-álgebra de tipo
finito, seja $\mathfrak q$ um ideal primo de $B$, e seja $\mathfrak p$
sua imagem inversa em $A$. Suponhamos que $\mathfrak q$ é ao mesmo tempo
maximal e minimal no conjunto de ideais primos de $B$ cuja imagem
inversa é $\mathfrak p$. Então existem um elemento
$g\in B-\mathfrak q$, uma $A$-álgebra finita $A'$, e um elemento
$f'\in A'$ tais que as $A$-álgebras $B_g$ e $A'_{f'}$ são
isomorfas.
\end{corollary}

\begin{proof}
Apliquemos \sref{III.4.4.5} a $Y=\Spec(A)$ e $X=\Spec(B)$. A
hipótese sobre $\mathfrak q$ diz precisamente que $\mathfrak q$
está isolado em sua fibra \sref[I]{I.1.1.7}.
\end{proof}

A consequência seguinte parece menos geral.

\oldpage[III]{137}

\begin{corollary}[4.4.7]
\label{III.4.4.7}
Seja $A$ um anel local noetheriano, seja $B$ uma $A$-álgebra de tipo
finito, e seja $\mathfrak n$ um ideal primo de $B$ cuja imagem inversa
em $A$ é o ideal maximal $\mathfrak m$. Suponhamos que $\mathfrak n$
é maximal em $B$ e minimal no conjunto de ideais primos de $B$
cuja imagem inversa é $\mathfrak m$ (o que significa também que
$\mathfrak mB$ é $\mathfrak n$-primário). Então existem uma
$A$-álgebra finita $A'$ e um ideal maximal $\mathfrak m'$ de $A'$
(cuja imagem inversa em $A$ é $\mathfrak m$) tais que $B_{\mathfrak
n}$ é isomorfo, como $A$-álgebra, a $A'_{\mathfrak m'}$.
\end{corollary}

O caso particular seguinte de \sref{III.4.4.7} também é chamado, às
vezes, de ``teorema principal''.

\begin{corollary}[4.4.8]
\label{III.4.4.8}
Sob as hipóteses de \sref{III.4.4.7}, suponhamos além disso que $A$ e
$B$ são domínios íntegros com o mesmo corpo de frações $K$.
Se $A$ é integralmente fechado, então $A=B$.
\end{corollary}

\begin{proof}
A Observação \sref{III.4.4.4} mostra que, ao aplicar
\sref{III.4.4.7}, podemos supor que $A'$ é um domínio íntegro com
corpo de frações $K$. A hipótese sobre $A$ dá então
$A'=A$, logo $B_{\mathfrak n}=A$. Como
$A\subset B\subset B_{\mathfrak n}$, segue-se que $A=B$.
\end{proof}

O enunciado \sref{III.4.4.8} é a forma na qual Zariski enunciou seu
``teorema principal'' (estendido aqui a anéis locais íntegros
noetherianos arbitrários).

Os corolários anteriores eram variantes locais de
\sref{III.4.4.3}, que é um resultado global. Eis outra
consequência de natureza global.

\begin{corollary}[4.4.9]
\label{III.4.4.9}
Seja $Y$ um esquema íntegro localmente noetheriano, e seja
$f:X\to Y$ um morfismo separado de tipo finito que é birracional.
Suponhamos além disso que $Y$ é normal e que toda fibra $f^{-1}(y)$,
para $y\in Y$, é finita. Então $f$ é uma imersão aberta. Se,
além disso, $f$ é fechado (em particular, se $f$ é próprio), então
$f$ é um isomorfismo.
\end{corollary}

\begin{proof}
Seja $x\in X$ e ponhamos $y=f(x)$. Como $f^{-1}(y)$ é um esquema
algébrico sobre $\kres(y)$, a hipótese de que é finito implica que
é discreto \sref[I]{I.6.4.4}. Além disso, $\sh{O}_y$ é integralmente
fechado, e $\sh{O}_x$ e $\sh{O}_y$ têm o mesmo corpo de frações
\sref[I]{I.7.1.5}. Podemos, portanto, aplicar
\sref{III.4.4.8}; se $f=(\psi,\theta)$, o homomorfismo
\[
  \theta_y^\#:\sh{O}_y\longrightarrow\sh{O}_x
\]
é bijetivo. De \sref[I]{I.6.5.4} segue-se que $f$ é um isomorfismo
local. Como $f$ é separado e $X$ é íntegro, $f$ é uma imersão
aberta \sref[I]{I.8.2.8}. A última afirmação se segue do fato de
que $f$ é dominante.
\end{proof}

\begin{proposition}[4.4.10]
\label{III.4.4.10}
Seja $Y$ um esquema localmente noetheriano, e seja $f:X\to Y$ um
morfismo localmente de tipo finito. O conjunto $X'$ de pontos
$x\in X$ que estão isolados em sua fibra $f^{-1}(f(x))$ é aberto em
$X$.
\end{proposition}

\begin{proof}
A questão é local sobre $X$ e $Y$, pelo que podemos supor que
$X$ e $Y$ são afins noetherianos e que $f$ é de tipo finito. O
morfismo $f$ é então afim e de tipo finito, logo
quase-projetivo \sref[II]{II.5.3.4}; basta aplicar
\sref{III.4.4.3}.
\end{proof}

\oldpage[III]{138}

\begin{corollary}[4.4.11]
\label{III.4.4.11}
Seja $Y$ um esquema localmente noetheriano, e seja $f:X\to Y$ um
morfismo próprio. O conjunto $U$ de pontos $y\in Y$ para os quais
$f^{-1}(y)$ é discreto é aberto em $Y$, e o morfismo restrito
$f^{-1}(U)\to U$ é finito. Em particular, todo morfismo próprio
quase-finito $X\to Y$ é finito.
\end{corollary}

\begin{proof}
O complementar de $U$ em $Y$ é a imagem por $f$ de $X-X'$, que
é fechado em $X$ por \sref{III.4.4.10}; como $f$ é uma aplicação
fechada, $U$ é aberto. Além disso, \sref[II]{II.6.2.2} mostra que
$f^{-1}(y)$ é finito para todo $y\in U$. Como o morfismo restrito
$f^{-1}(U)\to U$ é próprio \sref[II]{II.5.4.1}, é finito por
\sref{III.4.4.2}.
\end{proof}

\begin{remark}[4.4.12]
\label{III.4.4.12}
\emph{(i)}
Como se anunciou em \sref[II]{II.6.2.7}, demonstraremos no capítulo~V
que, se $Y$ é localmente noetheriano, todo morfismo separado
quase-finito $f:X\to Y$ é quase-afim, logo quase-projetivo. Seguir-se-á
que a conclusão do teorema principal \sref{III.4.4.3} continua
válida quando se supõe apenas que $f$ é separado e de tipo finito.
Com efeito, \sref{III.4.4.10} mostra que $X'$ é aberto em $X$; e
como $X$ é localmente noetheriano, a restrição de $f$ a $X'$ continua
sendo de tipo finito \sref[I]{I.6.3.5}, logo quase-finita pela
definição de $X'$, e claramente separada. Podemos, portanto, aplicar
\sref{III.4.4.3} a esta restrição.

\emph{(ii)}
No capítulo~IV daremos uma demonstração mais elementar de
\sref{III.4.4.10}, utilizando a teoria da dimensão.
\end{remark}


\subsection{Completamentos de módulos de homomorfismos}
\label{subsection:III.4.5}

\begin{proposition}[4.5.1]
\label{III.4.5.1}
Seja $A$ um anel noetheriano, seja $\mathfrak J$ um ideal de $A$, seja
$X$ um $A$-esquema de tipo finito, e sejam $\sh{F}$ e $\sh{G}$ dois
$\sh{O}_X$-módulos coerentes cujos suportes têm interseção
própria sobre $Y=\Spec(A)$ \sref[II]{II.5.4.10}. Então, para todo
inteiro $n\geq0$,
\[
  \Ext^n_{\sh{O}_X}(X;\sh{F},\sh{G})
\]
é um $A$-módulo finito, e seu completamento separado para a topologia
$\mathfrak J$-pré-ádica identifica-se canonicamente, com a notação de
\sref{III.4.1.7}, com
\[
  \Ext^n_{\sh{O}_{\hat{X}}}
    (\hat{X};\hat{\sh{F}},\hat{\sh{G}}).
\]
\end{proposition}

\begin{proof}
Sabemos por (T, 4.2) que existe uma sucessão espectral birregular
$\mathrm{E}(\sh{F},\sh{G})$ cujo termo limite é
\[
  \Ext_{\sh{O}_X}(X;\sh{F},\sh{G})
\]
e cujos termos $\mathrm{E}_2$ são
\[
  \mathrm{E}_2^{pq}
  =\HH^p\bigl(X,\shExt^q_{\sh{O}_X}(\sh{F},\sh{G})\bigr).
\]
O $\sh{O}_X$-módulo
$\shExt^q_{\sh{O}_X}(\sh{F},\sh{G})$ é coerente
\sref[0]{0.12.3.3}; seu suporte está contido na interseção dos
suportes de $\sh{F}$ e $\sh{G}$ (T, 4.2.2) e, consequentemente, é
próprio sobre $Y$ \sref[II]{II.5.4.10}. De \sref{III.3.2.4} se segue
que os $\mathrm{E}_2^{pq}$ são $A$-módulos finitos, e depois, de
\sref[0]{0.11.1.8}, que o mesmo ocorre com todo termo
$\mathrm{E}_r^{pq}$ da sucessão espectral e com seu termo limite.

Por outro lado, se $i:\hat{X}\to X$ é o morfismo canônico, então
$\hat{\sh{F}}$ e $\hat{\sh{G}}$ identificam-se canonicamente com
$i^*(\sh{F})$ e $i^*(\sh{G})$, e $i$ é plano
\sref[I]{I.10.8.8}, \sref[I]{I.10.8.9}. Para todo $q\geq0$ existe,
portanto, um $i$-morfismo canônico
\[
  u_q:\shExt^q_{\sh{O}_X}(\sh{F},\sh{G})
  \longrightarrow
  \shExt^q_{\sh{O}_{\hat{X}}}
    (\hat{\sh{F}},\hat{\sh{G}})
\]
\sref[0]{0.12.3.4}, e o $\sh{O}_{\hat{X}}$-homomorfismo
correspondente
\[
  u_q^\#:
  i^*\bigl(\shExt^q_{\sh{O}_X}(\sh{F},\sh{G})\bigr)
  \longrightarrow
  \shExt^q_{\sh{O}_{\hat{X}}}
    (\hat{\sh{F}},\hat{\sh{G}})
\]
é um isomorfismo \sref[0]{0.12.3.5}. Em outros termos,
\sref[I]{I.10.8.8} identifica
\[
  \shExt^q_{\sh{O}_{\hat{X}}}
    (\hat{\sh{F}},\hat{\sh{G}})
\]
canonicamente com o completamento de
$\shExt^q_{\sh{O}_X}(\sh{F},\sh{G})$, com a notação de
\sref{III.4.1.7}.

O teorema de comparação \sref{III.4.1.10} mostra agora, para todo
$p\geq0$, que
\[
  \HH^p\bigl(\hat{X},
    \shExt^q_{\sh{O}_{\hat{X}}}
      (\hat{\sh{F}},\hat{\sh{G}})\bigr)
\]
identifica-se canonicamente com o completamento separado, para a
topologia $\mathfrak J$-pré-ádica, de
\[
  \HH^p\bigl(X,\shExt^q_{\sh{O}_X}(\sh{F},\sh{G})\bigr).
\]

\oldpage[III]{139}

Denotemos por $\mathrm{E}(\hat{\sh{F}},\hat{\sh{G}})$ a sucessão
espectral birregular de (T, 4.2) associada a $\hat{\sh{F}}$ e
$\hat{\sh{G}}$. Se $\hat{A}$ é o completamento separado de $A$ para
a topologia $\mathfrak J$-pré-ádica, então, salvo isomorfismo
canônico,
\[
  \mathrm{E}_2^{pq}(\hat{\sh{F}},\hat{\sh{G}})
  =
  \mathrm{E}_2^{pq}(\sh{F},\sh{G})\otimes_A\hat{A}
\]
\sref[0]{0.7.3.3}.

O morfismo plano $i$ define um homomorfismo canônico de sucessões
espectrais
\[
  \varphi:\mathrm{E}(\sh{F},\sh{G})
  \longrightarrow
  \mathrm{E}(\hat{\sh{F}},\hat{\sh{G}})
  =
  \mathrm{E}(i^*(\sh{F}),i^*(\sh{G})).
\]
Sobre os termos $\mathrm{E}_2$ e sobre o termo limite,
respectivamente, este é o homomorfismo
\[
  \varphi_2^{pq}:
  \HH^p\bigl(X,\shExt^q_{\sh{O}_X}(\sh{F},\sh{G})\bigr)
  \longrightarrow
  \HH^p\bigl(\hat{X},
    \shExt^q_{\sh{O}_{\hat{X}}}
      (\hat{\sh{F}},\hat{\sh{G}})\bigr)
\]
e o homomorfismo
\[
  \varphi^n:
  \Ext^n_{\sh{O}_X}(X;\sh{F},\sh{G})
  \longrightarrow
  \Ext^n_{\sh{O}_{\hat{X}}}
    (\hat{X};\hat{\sh{F}},\hat{\sh{G}}),
\]
deduzidos de $u_q$ e $u_0$, respectivamente, por funtorialidade
\sref[0]{0.12.3.4}.

Após tensorizar com $\hat{A}$, os $\varphi_r^{pq}$ e
$\varphi^n$ dão homomorfismos de $\hat{A}$-módulos
\[
  \psi_r^{pq}:
  \mathrm{E}_r^{pq}(\sh{F},\sh{G})\otimes_A\hat{A}
  \longrightarrow
  \mathrm{E}_r^{pq}(\hat{\sh{F}},\hat{\sh{G}})
\]
e
\[
  \psi^n:
  \Ext^n_{\sh{O}_X}(X;\sh{F},\sh{G})\otimes_A\hat{A}
  \longrightarrow
  \Ext^n_{\sh{O}_{\hat{X}}}
    (\hat{X};\hat{\sh{F}},\hat{\sh{G}}).
\]
Como $\hat{A}$ é um $A$-módulo plano \sref[0]{0.7.3.3}, os
$\hat{A}$-módulos
\[
  \mathrm{E}_r^{pq}(\sh{F},\sh{G})\otimes_A\hat{A}
\]
formam uma sucessão espectral birregular cujo termo limite é
\[
  \Ext^n_{\sh{O}_X}(X;\sh{F},\sh{G})\otimes_A\hat{A},
\]
e os $\psi_r^{pq}$ e $\psi^n$ formam um morfismo de sucessões
espectrais. Como os $\psi_2^{pq}$ são isomorfismos, o mesmo ocorre
com os $\psi^n$ \sref[0]{0.11.1.5}.
\end{proof}

\begin{corollary}[4.5.2]
\label{III.4.5.2}
Sob as hipóteses de \sref{III.4.5.1}, suponhamos além disso que $A$ é
um anel noetheriano $\mathfrak J$-ádico. Então, para todo inteiro
$n\geq0$,
\[
  \Ext^n_{\sh{O}_X}(X;\sh{F},\sh{G})
  \isoto
  \Ext^n_{\sh{O}_{\hat{X}}}
    (\hat{X};\hat{\sh{F}},\hat{\sh{G}})
\]
canonicamente.
\end{corollary}

\begin{proof}
O $A$-módulo finito
$\Ext^n_{\sh{O}_X}(X;\sh{F},\sh{G})$ está separado e completo para a
topologia $\mathfrak J$-pré-ádica \sref[0]{0.7.3.6}.
\end{proof}

O caso particular $n=0$ de \sref{III.4.5.1} pode ser enunciado como
segue.

\begin{corollary}[4.5.3]
\label{III.4.5.3}
Sob as hipóteses de \sref{III.4.5.1}, para todo homomorfismo
$u:\sh{F}\to\sh{G}$, denotemos por
$\hat{u}:\hat{\sh{F}}\to\hat{\sh{G}}$ o homomorfismo completado
\sref[I]{I.10.8.4}. Existe um isomorfismo canônico
\[
\label{III.4.5.3.1}
  \left(\Hom_{\sh{O}_X}(\sh{F},\sh{G})\right)^\wedge
  \isoto
  \Hom_{\sh{O}_{\hat{X}}}(\hat{\sh{F}},\hat{\sh{G}}),
  \tag{4.5.3.1}
\]
onde o membro esquerdo é o completamento separado, para a
topologia $\mathfrak J$-pré-ádica, do $A$-módulo
$\Hom_{\sh{O}_X}(\sh{F},\sh{G})$. Este isomorfismo obtém-se passando
aos completamentos separados no homomorfismo
$u\mapsto\hat{u}$.
\end{corollary}


\subsection{Relações entre morfismos formais e morfismos ordinários}
\label{subsection:III.4.6}

\begin{proposition}[4.6.1]
\label{III.4.6.1}
Sejam $Y$ um esquema localmente noetheriano, $f:X\to Y$ um
morfismo próprio, $\sh{F}$ um $\sh{O}_X$-módulo coerente que é
plano sobre $Y$, e $y\in Y$. Suponhamos que, para algum inteiro
$n$,
\[
  \HH^n\bigl(f^{-1}(y),
    \sh{F}\otimes_{\sh{O}_Y}\kres(y)\bigr)=0.
\]
Então existe uma vizinhança $U$ de $y$ em $Y$ tal que
\[
  \RR^n f_*(\sh{F})|U=0,
\]
e, para todo inteiro $p\geq0$, o homomorfismo canônico
\[
  \bigl(\RR^{n-1}f_*(\sh{F})\bigr)_y
  \longrightarrow
  \HH^{n-1}\bigl(f^{-1}(y),
    \sh{F}\otimes_{\sh{O}_Y}
      (\sh{O}_y/\mathfrak m_y^{p+1})\bigr)
\]
\eqref{III.4.2.1.1} é sobrejetivo.
\end{proposition}

\begin{proof}
O $\sh{O}_Y$-módulo $\RR^n f_*(\sh{F})$ é coerente
\sref{III.3.2.1}; portanto, a primeira afirmação se seguirá uma vez que provemos
\[
  \bigl(\RR^n f_*(\sh{F})\bigr)_y=0
\]
\sref[0]{0.5.2.2}. Por \sref{III.4.2.1}, basta provar que
\[
  \HH^n\bigl(f^{-1}(y),
    \sh{F}\otimes_{\sh{O}_Y}
      (\sh{O}_y/\mathfrak m_y^{p+1})\bigr)=0
\]
para todo $p$. Isto é certo por hipótese para $p=0$; procedemos por
indução sobre $p$.

Ponhamos
\[
  X_p=X\times_Y
    \Spec(\sh{O}_y/\mathfrak m_y^{p+1}).
\]
Então $X_{p-1}$ é um subesquema fechado de $X_p$ com o mesmo
espaço subjacente \sref[I]{I.3.6.1}. A hipótese de indução
\[
  \HH^n\bigl(X_{p-1},
    \sh{F}\otimes_{\sh{O}_Y}
      (\sh{O}_y/\mathfrak m_y^p)\bigr)=0
\]
dá, portanto,
\[
  \HH^n\bigl(X_p,
    \sh{F}\otimes_{\sh{O}_Y}
      (\sh{O}_y/\mathfrak m_y^p)\bigr)=0.
\]
Por outro lado, a sucessão exata de $\sh{O}_{X_p}$-módulos
\[
  0\longrightarrow
  \mathfrak m_y^p\sh{F}/\mathfrak m_y^{p+1}\sh{F}
  \longrightarrow
  \sh{F}/\mathfrak m_y^{p+1}\sh{F}
  \longrightarrow
  \sh{F}/\mathfrak m_y^p\sh{F}
  \longrightarrow0
\]
dá a sucessão exata de cohomologia
\[
  \HH^n\bigl(X_p,
    \mathfrak m_y^p\sh{F}/\mathfrak m_y^{p+1}\sh{F}\bigr)
  \longrightarrow
  \HH^n\bigl(X_p,\sh{F}/\mathfrak m_y^{p+1}\sh{F}\bigr)
  \longrightarrow
  \HH^n\bigl(X_p,\sh{F}/\mathfrak m_y^p\sh{F}\bigr).
\]
Basta, pois, demonstrar que
\[
\label{III.4.6.1.1}
  \HH^n\bigl(X_p,
    \mathfrak m_y^p\sh{F}/\mathfrak m_y^{p+1}\sh{F}\bigr)=0.
  \tag{4.6.1.1}
\]
Com efeito, o termo central será então um submódulo do último termo,
que é zero pela hipótese de indução.

\oldpage[III]{140}

A fibra
\[
  Z=f^{-1}(y)=X\times_Y\Spec(\kres(y))
\]
é um subesquema fechado de $X_p$, e
$\mathfrak m_y^p\sh{F}/\mathfrak m_y^{p+1}\sh{F}$ é anulado por
$\mathfrak m_y$. Portanto, pode ser considerado um
$\sh{O}_Z$-módulo, e
\[
  \HH^n\bigl(Z,
    \mathfrak m_y^p\sh{F}/\mathfrak m_y^{p+1}\sh{F}\bigr)
  =
  \HH^n\bigl(X_p,
    \mathfrak m_y^p\sh{F}/\mathfrak m_y^{p+1}\sh{F}\bigr).
\]
Provaremos agora que o $\sh{O}_Z$-homomorfismo canônico
\[
\label{III.4.6.1.2}
  (\sh{F}/\mathfrak m_y\sh{F})
    \otimes_{\kres(y)}
      (\mathfrak m_y^p/\mathfrak m_y^{p+1})
  \longrightarrow
  \mathfrak m_y^p\sh{F}/\mathfrak m_y^{p+1}\sh{F}
  \tag{4.6.1.2}
\]
é bijetivo. Uma vez estabelecido isto, o fato de que
$\mathfrak m_y^p/\mathfrak m_y^{p+1}$ seja um
$\kres(y)$-módulo livre dá
\[
\begin{split}
  \HH^n\bigl(Z,
    \mathfrak m_y^p\sh{F}/\mathfrak m_y^{p+1}\sh{F}\bigr)
  &=
  \HH^n\bigl(Z,\sh{F}/\mathfrak m_y\sh{F}\bigr)
    \otimes_{\kres(y)}
      (\mathfrak m_y^p/\mathfrak m_y^{p+1})\\
  &=0
\end{split}
\]
\sref[0]{0.12.2.3}, dado que
$\HH^n(Z,\sh{F}/\mathfrak m_y\sh{F})=0$ por hipótese. Isto prova
\eqref{III.4.6.1.1}.

Para terminar a demonstração da primeira afirmação, resta demonstrar
que \eqref{III.4.6.1.2} é bijetivo. A questão pode ser verificada ponto a ponto em
$X$, e $\sh{F}_x$ é um $\sh{O}_y$-módulo plano por hipótese para
todo $x\in f^{-1}(y)$. Basta, pois, aplicar
\sref[0]{0.10.2.1}, c), dado que
$\sh{F}_x/\mathfrak m_y\sh{F}_x$ é plano sobre o corpo
$\kres(y)=\sh{O}_y/\mathfrak m_y$.

Para a segunda afirmação de \sref{III.4.6.1}, reduzimo-nos imediatamente, como
em \sref{III.4.2.1}, ao caso em que $Y$ é afim e $y$ é
fechado. O mesmo argumento demonstra, para todo $k>0$, que
\[
\label{III.4.6.1.3}
  \HH^n\bigl(X_{k+1},
    \mathfrak m_y^k\sh{F}/
      \mathfrak m_y^{k+p+1}\sh{F}\bigr)=0.
  \tag{4.6.1.3}
\]

\oldpage[III]{141}

De \sref{III.4.2.1} segue-se que
\[
\label{III.4.6.1.4}
  \bigl(\RR^n f_*(\mathfrak m_y^k\sh{F})\bigr)_y=0.
  \tag{4.6.1.4}
\]
A sucessão exata de cohomologia dá agora a exatidão de
\[
  \bigl(\RR^{n-1}f_*(\sh{F})\bigr)_y
  \longrightarrow
  \bigl(\RR^{n-1}f_*
    (\sh{F}/\mathfrak m_y^p\sh{F})\bigr)_y
  \longrightarrow
  \bigl(\RR^n f_*(\mathfrak m_y^p\sh{F})\bigr)_y=0.
\]
Como $y$ é fechado e $Y$ é afim, temos
\sref{III.1.4.11}
\[
  \RR^{n-1}f_*(\sh{F}/\mathfrak m_y^p\sh{F})
  =
  \bigl(\HH^{n-1}
    (X,\sh{F}/\mathfrak m_y^p\sh{F})\bigr)^{\sim}
  =
  \bigl(\HH^{n-1}
    (f^{-1}(y),\sh{F}/\mathfrak m_y^p\sh{F})\bigr)^{\sim}
\]
(G, II, 4.9.1). O módulo
\[
  \HH^{n-1}
    (f^{-1}(y),\sh{F}/\mathfrak m_y^p\sh{F})
\]
é um $\sh{O}_y/\mathfrak m_y^p$-módulo, e portanto
\[
  \bigl(\RR^{n-1}f_*
    (\sh{F}/\mathfrak m_y^p\sh{F})\bigr)_y
  =
  \HH^{n-1}
    (f^{-1}(y),\sh{F}/\mathfrak m_y^p\sh{F}).
\]
Isto termina a demonstração.
\end{proof}

\begin{corollary}[4.6.2]
\label{III.4.6.2}
Sejam $Y$ um esquema localmente noetheriano, $f:X\to Y$ um
morfismo próprio e plano, $\sh{F}$ e $\sh{G}$ dois
$\sh{O}_X$-módulos localmente livres, e $y$ um ponto de $Y$. Ponhamos
\[
  X_y=f^{-1}(y)=X\times_Y\Spec(\kres(y)),\qquad
  \sh{F}_y=\sh{F}\otimes_{\sh{O}_Y}\kres(y),\qquad
  \sh{G}_y=\sh{G}\otimes_{\sh{O}_Y}\kres(y),
\]
e suponhamos que
\[
\label{III.4.6.2.1}
  \HH^1\bigl(
    X_y,\shHom_{\sh{O}_{X_y}}(\sh{F}_y,\sh{G}_y)
  \bigr)=0.
  \tag{4.6.2.1}
\]
Então, para todo homomorfismo $u_0:\sh{F}_y\to\sh{G}_y$, existem
uma vizinhança aberta $U$ de $y$ em $Y$ e um homomorfismo
\[
  u:\sh{F}|_{f^{-1}(U)}\longrightarrow
    \sh{G}|_{f^{-1}(U)}
\]
tal que $u_0=u\otimes1$.
\end{corollary}

\begin{proof}
A hipótese nos permite aplicar \sref{III.4.6.1}, com $n=1$
e $p=0$, ao $\sh{O}_X$-módulo coerente
\[
  \sh{H}=\shHom_{\sh{O}_X}(\sh{F},\sh{G}).
\]
Com efeito, $\sh{H}$ é localmente livre e, portanto, com maior razão, plano sobre
$Y$, enquanto que o $\sh{O}_{X_y}$-módulo
$\sh{H}\otimes_{\sh{O}_Y}\kres(y)$ identifica-se com
$\shHom_{\sh{O}_{X_y}}(\sh{F}_y,\sh{G}_y)$
\sref[0]{0.6.2.2}. Podemos supor que $Y=\Spec(A)$ é afim.
Então \sref{III.1.4.11} dá
\[
  \RR^0f_*(\sh{H})
  =
  \bigl(\Hom_{\sh{O}_X}(\sh{F},\sh{G})\bigr)^{\sim},
\]
e, consequentemente,
\[
  \bigl(\RR^0f_*(\sh{H})\bigr)_y
  =
  \Hom_{\sh{O}_X}(\sh{F},\sh{G})
    \otimes_A\sh{O}_y.
\]
Por \sref{III.4.6.1}, o homomorfismo canônico
\[
  \Hom_{\sh{O}_X}(\sh{F},\sh{G})
    \otimes_A\sh{O}_y
  \longrightarrow
  \Hom_{\sh{O}_{X_y}}(\sh{F}_y,\sh{G}_y)
\]
é sobrejetivo. Isto prova o corolário, dado que todo elemento do
módulo da esquerda pode escrever-se como $u\otimes(1/s)$, onde
$s\in A\setminus\mathfrak m_y$.
\end{proof}

\begin{corollary}[4.6.3]
\label{III.4.6.3}
Sob as hipóteses de \sref{III.4.6.2}, se $u_0$ é injetivo
(respectivamente sobrejetivo, bijetivo), pode-se escolher $u$
injetivo (respectivamente sobrejetivo, bijetivo).
\end{corollary}

\begin{proof}
Podemos restringir-nos ao caso $U=Y$. Basta provar
que, se $u_0$ é injetivo (respectivamente sobrejetivo), então
\[
  \Ker u_x=0
  \quad\text{(respectivamente }\Coker u_x=0\text{)}
\]
para todo $x\in f^{-1}(y)$. Com efeito, $\Ker u$ e $\Coker u$ são
$\sh{O}_X$-módulos coerentes \sref[0]{0.5.3.4}; portanto, existe uma
vizinhança $V$ de $f^{-1}(y)$ em $X$ sobre a qual a restrição de
$\Ker u$ (respectivamente $\Coker u$) é zero \sref[0]{0.5.2.2}.
Como $f$ é fechado, existe uma vizinhança aberta
$U'\subset U$ de $y$ tal que $f^{-1}(U')\subset V$, o que prova
a afirmação depois de substituir $U$ por $U'$.

\oldpage[III]{142}

Por hipótese,
\[
  u_x\otimes1:
  \sh{F}_x\otimes_{\sh{O}_y}\kres(y)
  \longrightarrow
  \sh{G}_x\otimes_{\sh{O}_y}\kres(y)
\]
é injetivo (respectivamente sobrejetivo), $\sh{F}_x$ e $\sh{G}_x$
são $\sh{O}_x$-módulos livres finitos, e $\sh{O}_x$ é plano sobre
$\sh{O}_y$. Se $u_x\otimes1$ é injetivo, a injetividade de $u_x$
se segue de \sref[0]{0.10.2.4}. Se $u_x\otimes1$ é sobrejetivo,
então, com maior razão, o homomorfismo induzido
\[
  \sh{F}_x/\mathfrak m_x\sh{F}_x
  \longrightarrow
  \sh{G}_x/\mathfrak m_x\sh{G}_x
\]
é sobrejetivo. Como $\sh{G}_x$ é um $\sh{O}_x$-módulo finito e
$\sh{O}_x$ é um anel local de ideal maximal $\mathfrak m_x$, a
conclusão se segue do lema de Nakayama (Bourbaki,
\emph{Alg\`ebre}, cap.~VIII, \S\,6, n.º~3, cor.~4 da prop.~6).
\end{proof}

\begin{corollary}[4.6.4]
\label{III.4.6.4}
Sejam $Y$ um esquema localmente noetheriano, $f:X\to Y$ um
morfismo próprio e plano, e $y$ um ponto de $Y$, e ponhamos
$X_y=X\times_Y\Spec(\kres(y))$. Seja $\sh{E}_0$ um
$\sh{O}_{X_y}$-módulo localmente livre tal que
\[
\label{III.4.6.4.1}
  \HH^1\bigl(
    X_y,\shHom_{\sh{O}_{X_y}}(\sh{E}_0,\sh{E}_0)
  \bigr)=0.
  \tag{4.6.4.1}
\]
Sejam $\sh{F}$ e $\sh{G}$ dois $\sh{O}_X$-módulos localmente livres
tais que $\sh{F}_y$ e $\sh{G}_y$, com a notação de
\sref{III.4.6.2}, são isomorfos a $\sh{E}_0$. Então existe uma
vizinhança aberta $U$ de $y$ tal que
$\sh{F}|_{f^{-1}(U)}$ e $\sh{G}|_{f^{-1}(U)}$ são isomorfos.
\end{corollary}

\begin{corollary}[4.6.5]
\label{III.4.6.5}
Sob as hipóteses de \sref{III.4.6.4} sobre $f$, $X$ e $Y$,
suponhamos que
\[
  \HH^1(X_y,\sh{O}_{X_y})=0.
\]
Se $\sh{F}$ e $\sh{G}$ são dois $\sh{O}_X$-módulos invertíveis tais
que $\sh{F}_y$ e $\sh{G}_y$ são isomorfos, então existe uma vizinhança
aberta $U$ de $y$ tal que $\sh{F}|_{f^{-1}(U)}$ e
$\sh{G}|_{f^{-1}(U)}$ são isomorfos.
\end{corollary}

\begin{proof}
Basta aplicar \sref{III.4.6.4} aos módulos
$\sh{F}^{-1}\otimes\sh{G}$ e $\sh{O}_X$.
\end{proof}

\begin{remarks}[4.6.6]
\label{III.4.6.6}
\begin{enumerate}
\item[(i)] Utilizando \sref{III.4.6.5}, estabeleceremos no capítulo~V
  a classificação dos feixes invertíveis sobre um fibrado projetivo
  anunciada em \sref[II]{II.4.2.7}.
\item[(ii)] O resultado de \sref{III.4.6.1} aparecerá em \S\,7 como
  consequência de proposições mais gerais.
\end{enumerate}
\end{remarks}

\begin{proposition}[4.6.7]
\label{III.4.6.7}
Seja $Z$ um esquema localmente noetheriano, e sejam $X$ e $Y$
dois $Z$-esquemas cujos morfismos estruturais
\[
  g:X\longrightarrow Z,\qquad h:Y\longrightarrow Z
\]
são próprios. Seja $f:X\to Y$ um $Z$-morfismo, seja $z$ um ponto
de $Z$, e ponhamos
\[
  f_z=f\times_Z1:
  X\times_Z\Spec(\kres(z))
  \longrightarrow
  Y\times_Z\Spec(\kres(z)).
\]
\begin{enumerate}
\item[(i)] Se $f_z$ é finito (respectivamente uma imersão fechada),
  existe uma vizinhança aberta $U$ de $z$ tal que o morfismo
  \[
    g^{-1}(U)\longrightarrow h^{-1}(U)
  \]
  induzido por $f$ é finito (respectivamente uma imersão fechada).
\item[(ii)] Suponhamos além disso que $g$ é plano. Se $f_z$ é um
  isomorfismo, existe uma vizinhança aberta $U$ de $z$ tal que
  o morfismo $g^{-1}(U)\to h^{-1}(U)$ induzido por $f$ é um
  isomorfismo.
\end{enumerate}
\end{proposition}

\begin{proof}
Em ambos os casos basta provar que, para todo
$y\in h^{-1}(z)$, existe uma vizinhança $V_y$ de $y$ tal que
a restrição
\[
  f^{-1}(V_y)\longrightarrow V_y
\]
de $f$ é finita (respectivamente uma imersão fechada, um isomorfismo).
Com efeito, se $V$ é a união dos $V_y$, então
$f^{-1}(V)\to V$ é finito (respectivamente uma imersão fechada, um
isomorfismo) \sref[II]{II.6.1.1}, \sref[I]{I.4.2.4}. Como $h$ é
fechado, existe uma vizinhança aberta $U$ de $z$ tal que
$h^{-1}(U)\subset V$, e a proposição se segue.

\noindent\textit{(i)} Observe-se primeiro que $f$ é próprio
\sref[II]{II.5.4.3}. Se
$f_z$ é finito, a vizinhança requerida $V_y$, para cada
$y\in h^{-1}(z)$, existe por \sref{III.4.4.11}.

\oldpage[III]{143}

Para tratar o caso em que $f_z$ é uma imersão fechada, podemos
supor já, portanto, que $f$ é finito. Assim,
$X=\Spec(\sh{B})$, onde $\sh{B}$ é uma
$\sh{O}_Y$-álgebra coerente, e $f$ corresponde \sref[II]{II.1.2.7} ao
homomorfismo canônico
\[
  u:\sh{O}_Y\longrightarrow\sh{B}.
\]
Se, para todo $y\in h^{-1}(z)$, o homomorfismo
$u_y:\sh{O}_y\to\sh{B}_y$ é sobrejetivo, então $u$ é sobrejetivo em
alguma vizinhança $V_y$ de $y$, dado que $\Coker u$ é coerente
\sref[0]{0.5.3.4}, \sref[0]{0.5.2.2}. Agora o morfismo finito
$f_z$ corresponde a
\[
  v=u\otimes1:
  \sh{O}_Y\otimes_{\sh{O}_Z}\kres(z)
  \longrightarrow
  \sh{B}\otimes_{\sh{O}_Z}\kres(z),
\]
e a hipótese de que $f_z$ é uma imersão fechada implica que
\[
  u_y\otimes1:
  \sh{O}_y\otimes_{\sh{O}_z}\kres(z)
  \longrightarrow
  \sh{B}_y\otimes_{\sh{O}_z}\kres(z)
\]
é sobrejetivo. Como $\sh{B}_y$ é um $\sh{O}_y$-módulo finito e
$\sh{O}_y$ é um anel local noetheriano, a conclusão se segue, como
em \sref{III.4.6.3}, do lema de Nakayama.

\noindent\textit{(ii)} O mesmo raciocínio mostra que agora
basta provar
que $u_y$ é bijetivo, sabendo que $u_y\otimes1$ é bijetivo.
A injetividade requerida se segue do lema seguinte; a sobrejetividade
se segue do argumento precedente.
\end{proof}

\begin{lemma}[4.6.7.1]
\label{III.4.6.7.1}
Sejam $A$ e $B$ anéis locais noetherianos, seja
$\rho:A\to B$ um homomorfismo local, e seja $u:N\to M$ um
homomorfismo de $B$-módulos. Suponhamos que $M$ é plano sobre $A$,
que $N$ é um $B$-módulo finito, e que
\[
  u\otimes1:
  N\otimes_A k\longrightarrow M\otimes_A k,
\]
onde $k$ é o corpo residual de $A$, é injetivo. Então $N$ é
plano sobre $A$ e $u$ é injetivo.
\end{lemma}

\begin{proof}
Para provar a primeira afirmação, devemos demonstrar que, para todo par de
$A$-módulos finitos $P,Q$ e todo $A$-homomorfismo injetivo
$v:P\to Q$, o homomorfismo
\[
  1_N\otimes v:N\otimes_A P\longrightarrow N\otimes_A Q
\]
é injetivo. Consideremos o diagrama comutativo
\[
  \xymatrix{
    N\otimes_A P
      \ar[r]^-{1_N\otimes v}
      \ar[d]_{u\otimes1_P}
      &
    N\otimes_A Q
      \ar[d]^{u\otimes1_Q}
    \\
    M\otimes_A P
      \ar[r]_{1_M\otimes v}
      &
    M\otimes_A Q .
  }
\]
Como $1_M\otimes v$ é injetivo pela planicidade de $M$, basta
provar que $u\otimes1_P$ é injetivo.

Seja $\mathfrak m$ o ideal maximal de $A$. A
filtração $\mathfrak m$-ádica sobre o $A$-módulo $N\otimes_A P$ é
também sua filtração $\mathfrak mB$-ádica como $B$-módulo. A
topologia definida por esta filtração é separada: $B$ é
noetheriano, $\mathfrak mB$ está contido no radical de $B$, e
$N\otimes_A P$ é um $B$-módulo finito, dado que $N$ é finito sobre $B$
e $P$ é finito sobre $A$ \sref[0]{0.7.3.5}. Basta, portanto,
provar que
\[
  \operatorname{gr}(u\otimes1_P):
  \operatorname{gr}_\bullet(N\otimes_A P)
  \longrightarrow
  \operatorname{gr}_\bullet(M\otimes_A P)
\]
é injetivo, onde os módulos graduados associados se tomam com
respeito às filtrações $\mathfrak m$-ádicas (Bourbaki,
\emph{Alg\`ebre commutative}, cap.~III, \S\,2, n.º~8, cor.~1 do
teor.~1).

Como $M$ é plano sobre $A$, os homomorfismos
\[
  M\otimes_A(\mathfrak m^nP)
  \longrightarrow
  \mathfrak m^n(M\otimes_A P)
\]
são bijetivos. Consequentemente, o homomorfismo canônico
\[
  \vphi_M:
  \operatorname{gr}_0(M)\otimes_A\operatorname{gr}_\bullet(P)
  \longrightarrow
  \operatorname{gr}_\bullet(M\otimes_A P)
\]
é bijetivo.

\oldpage[III]{144}

Temos o diagrama comutativo
\[
  \xymatrix{
    \operatorname{gr}_0(N)\otimes_A\operatorname{gr}_\bullet(P)
      \ar[r]^-{\operatorname{gr}_0(u)\otimes1}
      \ar[d]_{\vphi_N}
      &
    \operatorname{gr}_0(M)\otimes_A\operatorname{gr}_\bullet(P)
      \ar[d]^{\vphi_M}
    \\
    \operatorname{gr}_\bullet(N\otimes_A P)
      \ar[r]_{\operatorname{gr}(u\otimes1)}
      &
    \operatorname{gr}_\bullet(M\otimes_A P).
  }
\]
Aqui $\vphi_M$ é bijetivo e $\vphi_N$ é sobrejetivo.
Além disso, $\operatorname{gr}_0(u)$ é injetivo por hipótese, e
\[
\begin{split}
  \operatorname{gr}_0(N)\otimes_A\operatorname{gr}_\bullet(P)
    &=
  \operatorname{gr}_0(N)\otimes_k\operatorname{gr}_\bullet(P),\\
  \operatorname{gr}_0(M)\otimes_A\operatorname{gr}_\bullet(P)
    &=
  \operatorname{gr}_0(M)\otimes_k\operatorname{gr}_\bullet(P).
\end{split}
\]
Assim, $\operatorname{gr}_0(u)\otimes1$ também é injetivo. Do
diagrama segue-se que $\operatorname{gr}(u\otimes1_P)$ é
injetivo, o que prova a primeira afirmação. A segunda se segue do
mesmo argumento tomando $P=A$.
\end{proof}

\begin{proposition}[4.6.8]
\label{III.4.6.8}
Seja $Z$ um esquema localmente noetheriano, e sejam $X$ e $Y$
dois $Z$-esquemas cujos morfismos estruturais
$g:X\to Z$ e $h:Y\to Z$ são próprios. Seja $Z'$ um subconjunto fechado
de $Z$, ponhamos
\[
  X'=g^{-1}(Z'),\qquad Y'=h^{-1}(Z'),
\]
e denotemos por
\[
  \widehat X=X_{/X'},\qquad
  \widehat Y=Y_{/Y'},\qquad
  \widehat Z=Z_{/Z'}
\]
os completamentos formais ao longo destes subconjuntos fechados. Seja $f:X\to Y$
um $Z$-morfismo e seja
$\widehat f:\widehat X\to\widehat Y$ sua extensão aos completamentos
formais. Para que $\widehat f$ seja um isomorfismo
(respectivamente uma imersão fechada), é necessário e suficiente
que exista uma vizinhança aberta $U$ de $Z'$ tal que o
morfismo
\[
  g^{-1}(U)\longrightarrow h^{-1}(U)
\]
induzido por $f$ seja um isomorfismo (respectivamente uma imersão fechada).
\end{proposition}

\begin{proof}
A suficiência da condição é imediata
\sref[I]{I.10.14.7}. Para provar sua necessidade, basta, pelo
mesmo raciocínio de \sref{III.4.6.7}, demonstrar que para todo
$y\in Y'$ existe uma vizinhança aberta $V_y$ de $y$ tal que
$f^{-1}(V_y)\to V_y$ seja um isomorfismo (respectivamente uma
imersão fechada). Ficamos assim reduzidos ao caso
\[
  Y=Z=\Spec(A),
\]
onde $A$ é noetheriano.

Por hipótese \sref[I]{I.10.9.1}, \sref[I]{I.10.14.2}, a fibra
$f^{-1}(y)$ consiste de um único ponto para todo $y\in Y'$. Como
$f$ é próprio \sref[II]{II.5.4.3}, \sref{III.4.4.11} dá uma vizinhança
aberta $U$ de $y$ tal que $f^{-1}(U)\to U$ é finito. Podemos,
portanto, supor que $f$ é finito, de modo que
\[
  X=\Spec(B),
\]
onde $B$ é uma $A$-álgebra finita.

Se $Y'=V(\mathfrak J)$, então
\[
  \widehat Y=\Spf(\widehat A),
  \qquad
  \widehat X=\Spf(\widehat B),
\]
onde $\widehat A$ é o completamento separado de $A$ para a
topologia $\mathfrak J$-ádica, e $\widehat B$ é o completamento
separado de $B$ para a topologia $\mathfrak JB$-ádica, ou,
equivalentemente, o completamento separado do $A$-módulo $B$ para
a topologia $\mathfrak J$-ádica. Além disso, $\widehat f$
corresponde à extensão contínua
\[
  \widehat{\vphi}:\widehat A\longrightarrow\widehat B
\]
do homomorfismo canônico de anéis $\vphi:A\to B$, e a
hipótese é que $\widehat{\vphi}$ é sobrejetivo (respectivamente
bijetivo) \sref[I]{I.10.14.2}. Mas $\widehat{\vphi}$ é também a
extensão contínua de $\vphi$ considerado como um homomorfismo de
$A$-módulos. De \sref[I]{I.10.8.14} segue-se que existe uma
vizinhança aberta $U$ de $Y'$ tal que a restrição a $U$ do
homomorfismo de $\sh{O}_Y$-módulos
\[
  \widetilde{\vphi}:\widetilde A\longrightarrow\widetilde B
\]
é sobrejetiva (respectivamente bijetiva), o que termina a demonstração.
\end{proof}


\subsection{Um critério de amplitude}
\label{subsection:III.4.7}

\begin{theorem}[4.7.1]
\label{III.4.7.1}
Sejam $Y$ um esquema localmente noetheriano, $f:X\to Y$ um
morfismo próprio, $\sh{L}$ um $\sh{O}_X$-módulo invertível,
e $y$ um ponto de $Y$. Ponhamos
\[
  X_y=X\times_Y\Spec(\kres(y))=f^{-1}(y),
\]
e seja $g:X_y\to X$ o morfismo canônico. Se
\[
  \sh{L}_y=g^*(\sh{L})
  =\sh{L}\otimes_{\sh{O}_Y}\kres(y)
\]
é amplo sobre $X_y$, então existe uma vizinhança aberta $U$ de $y$
em $Y$ tal que $\sh{L}|_{f^{-1}(U)}$ é amplo para a restrição
de $f$ a $f^{-1}(U)$.
\end{theorem}

\begin{proof}
\oldpage[III]{145}

\noindent\textit{I)} Ponhamos
\[
  Y'=\Spec(\sh{O}_y),\qquad
  X'=X\times_Y Y',\qquad
  \sh{L}'=\sh{L}\otimes_{\sh{O}_Y}\sh{O}_y.
\]
Provamos primeiro que $\sh{L}'$ é amplo para $f'=f_{(Y')}$. Temos
o diagrama comutativo
\[
  \xymatrix{
    X\ar[d]_f
      &
    X'\ar[l]\ar[d]_{f'}
      &
    X_y\ar[l]\ar[d]_{f_y}
    \\
    Y
      &
    Y'\ar[l]
      &
    \Spec(\kres(y))\ar[l].
  }
\]
Como $f'$ é próprio \sref[II]{II.5.4.2}, (iii), e $\sh{O}_y$ é
noetheriano, \sref[II]{II.4.6.6} mostra que podemos reduzir-nos ao caso
\[
  Y=Y'=\Spec(\sh{O}_y),\qquad X=X',
\]
supor que $\sh{L}_y$ é amplo para $f_y$, e provar que $\sh{L}$
é amplo para $f$. Aplicamos o critério
\sref{III.2.6.1}, c), e provamos que para todo
$\sh{O}_X$-módulo coerente $\sh{F}$ existe um inteiro $N$ tal que
\[
  \HH^1(X,\sh{F}(n))=0
  \quad\text{para todo }n\geq N,
  \qquad
  \sh{F}(n)=\sh{F}\otimes\sh{L}^{\otimes n}.
\]

O ponto $y$ é fechado em $Y$ e corresponde ao ideal maximal
$\mathfrak m$ de $\sh{O}_y$. Portanto, $X_y$ é o
subesquema fechado de $X$ definido pelo ideal coerente
\[
  \sh{J}
  =
  f^*(\widetilde{\mathfrak m})\sh{O}_X
  =
  \mathfrak m\sh{O}_X
\]
\sref[I]{I.4.4.5}, e $g$ é sua injeção canônica. Consideremos
a $\kres(y)$-álgebra graduada
\[
  S=\operatorname{gr}(\sh{O}_y)
   =\bigoplus_{j\geq0}
      \mathfrak m^j/\mathfrak m^{j+1}.
\]
É de tipo finito porque $\sh{O}_y$ é noetheriano. A
$\sh{O}_X$-álgebra
\[
  \sh{S}=f^*(\widetilde S)
\]
é quase-coerente e de tipo finito, e é anulada por
$\sh{J}$. Assim, se $\sh{S}_y=g^*(\sh{S})$, então $\sh{S}_y$ é uma
$\sh{O}_{X_y}$-álgebra quase-coerente de tipo finito e
\[
  \sh{S}=g_*(\sh{S}_y).
\]

Ponhamos
\[
  \sh{M}_j
  =
  \mathfrak m^j\sh{F}/\mathfrak m^{j+1}\sh{F},
  \qquad
  \sh{M}
  =
  \bigoplus_{j\geq0}\sh{M}_j
  =
  \operatorname{gr}(\sh{F}).
\]
Como $\sh{F}$ é coerente, $\sh{M}$ é um
$\sh{S}$-módulo quase-coerente de tipo finito \sref[0]{0.10.1.1}; também é
anulado por $\sh{J}$. Consequentemente, se
\[
  \sh{M}'_j=g^*(\sh{M}_j),
  \qquad
  \sh{M}'=g^*(\sh{M})
    =\bigoplus_{j\geq0}\sh{M}'_j,
\]
então $\sh{M}'$ é um $\sh{S}_y$-módulo graduado quase-coerente de
tipo finito e $\sh{M}=g_*(\sh{M}')$. Além disso, pondo
\[
  \sh{M}'_j(n)
  =
  \sh{M}'_j\otimes\sh{L}_y^{\otimes n},
\]
temos $\sh{M}'_j(n)=g^*(\sh{M}_j(n))$.

O morfismo $f_y$ é próprio \sref[II]{II.5.4.2}, (iii), e
$\sh{L}_y$ é amplo; portanto, $f_y$ é projetivo
\sref[II]{II.5.5.4}, \sref[II]{II.4.6.11}. Podemos, pois, aplicar
\sref{III.2.4.1}, (ii), a
$\Spec(\kres(y))$, $f_y$, $\sh{S}_y$, $\sh{L}_y$ e $\sh{M}'$.
Existe um inteiro $N$ tal que, para $n\geq N$,
\[
  \HH^q(X_y,\sh{M}'_j(n))=0
\]
para todo $q>0$ e todo $j$. Segue-se que
\[
  \HH^q(X,\sh{M}_j(n))=0
\]
para todo $q>0$ e todo $j$ (G, II, 4.9.1).

Ponhamos agora
\[
  \sh{F}(n)_j
  =
  \sh{F}(n)/\mathfrak m^{j+1}\sh{F}(n).
\]
Para $j\geq1$ temos
\[
  \sh{F}(n)_{j-1}
  =
  \sh{F}(n)_j/\sh{M}_j(n),
  \qquad
  \sh{F}(n)_0=\sh{M}_0(n).
\]
Assim, $\HH^1(X,\sh{F}(n)_0)=0$, e a sucessão exata de cohomologia
dá, para todo $j\geq1$,
\[
  \HH^1(X,\sh{F}(n)_j)
  =
  \HH^1(X,\sh{F}(n)_{j-1}).
\]
Portanto, $\HH^1(X,\sh{F}(n)_j)=0$ para todo $j\geq0$. Por
\sref{III.4.2.1}, segue-se que
\[
  \HH^1(X,\sh{F}(n))=0,
\]
o que prova a afirmação.

\oldpage[III]{146}

\noindent\textit{II)} Voltemos à notação do início da
demonstração. Podemos supor que $Y=\Spec(A)$ é afim. Como
$f'$ é de tipo finito e $\sh{L}'$ é amplo para $f'$, existe um
inteiro $m>0$ tal que $\sh{L}'^{\otimes m}$ é muito amplo para
$f'$ \sref[II]{II.4.6.11}. Substituindo $\sh{L}$ por
$\sh{L}^{\otimes m}$, basta considerar o caso em que
$\sh{L}'$ é muito amplo para $f'$ e provar que
$\sh{L}|_{f^{-1}(U)}$ é muito amplo para a restrição de $f$.

Como $f'$ é próprio, existe, para um inteiro adequado $r>0$, uma
$Y'$-imersão fechada
\[
  j:X'\longrightarrow P=\bb{P}_{Y'}^r
\]
tal que $\sh{L}'\simeq j^*(\sh{O}_P(1))$
\sref[II]{II.5.5.4}, (ii). Essa imersão fechada corresponde
canonicamente a um $\sh{O}_{X'}$-homomorfismo sobrejetivo
\[
  u:\sh{O}_{X'}^{r+1}\longrightarrow\sh{L}'
\]
\sref[II]{II.4.2.3}, ou, equivalentemente, \sref[0]{0.5.1.1} a
$r+1$ seções $s'_i$ $(0\leq i\leq r)$ de $\sh{L}'$ sobre $X'$
que geram este $\sh{O}_{X'}$-módulo.

Por definição, estas são também seções de $f'_*(\sh{L}')$ sobre
$Y'$, e
\[
  f'_*(\sh{L}')
  =
  f_*(\sh{L})\otimes_{\sh{O}_Y}\sh{O}_{Y'}
\]
\sref[0]{0.5.4.10}. Como $Y$ é afim e $\sh{O}_y$ é o
anel local de $A$ no ideal primo $\mathfrak j_y$, cada $s'_i$
tem a forma
\[
  s'_i=s''_i/t_i,
\]
onde $s''_i$ é uma seção de $f_*(\sh{L})$ sobre $Y$ e
$t_i\in A\setminus\mathfrak j_y$. Segue-se que existem uma
vizinhança aberta $V$ de $y$ em $Y$ e seções $s_i$ de
$f_*(\sh{L})|_V$ tais que $s'_i=s_i/1$ (o espaço $Y'$ está
contido em $V$; veja-se \sref[I]{I.2.4.2}).

As seções $s_i$ são seções de $\sh{L}$ sobre $f^{-1}(V)$ e
definem, portanto, um homomorfismo
\[
  v:
  \bigl(\sh{O}_X|_{f^{-1}(V)}\bigr)^{r+1}
  \longrightarrow
  \sh{L}|_{f^{-1}(V)}.
\]
Por construção, é sobrejetivo em todo ponto de $f^{-1}(y)$.
Como $\Coker(v)$ é coerente \sref[0]{0.5.3.4}, seu suporte é
fechado \sref[0]{0.5.2.2}; portanto, existe uma vizinhança aberta
$W\subset f^{-1}(V)$ de $f^{-1}(y)$ sobre a qual $v$ é sobrejetivo.
Como $f$ é fechado, após restringir em torno de $y$, podemos escolher uma
vizinhança aberta $U\subset V$ com $f^{-1}(U)\subset W$.
A conclusão se segue agora de \sref[II]{II.4.2.3}.
\end{proof}


\subsection{Morfismos finitos de esquemas formais}
\label{subsection:III.4.8}

\begin{proposition}[4.8.1]
\label{III.4.8.1}
Seja $\mathfrak{Y}$ um esquema formal localmente noetheriano,
$\sh{K}$ um ideal de definição de $\mathfrak{Y}$, e
$f:\mathfrak{X}\to\mathfrak{Y}$ um morfismo de esquemas formais.
As condições seguintes são equivalentes:
\begin{enumerate}
  \item[\rm{a)}] $\mathfrak{X}$ é localmente noetheriano, $f$ é um
    morfismo ádico \sref[I]{I.10.12.1}, e, se
    $\sh{J}=f^*(\sh{K})\sh{O}_{\mathfrak{X}}$, o morfismo
    \[
      f_0:
      (\mathfrak{X},\sh{O}_{\mathfrak{X}}/\sh{J})
      \longrightarrow
      (\mathfrak{Y},\sh{O}_{\mathfrak{Y}}/\sh{K})
    \]
    induzido por $f$ é finito.
  \item[\rm{b)}] $\mathfrak{X}$ é localmente noetheriano e é o
    limite indutivo de um sistema indutivo ádico de $(Y_n)$, $(X_n)$,
    tal que o morfismo $X_0\to Y_0$ é finito.
  \item[\rm{c)}] Todo ponto de $\mathfrak{Y}$ possui uma vizinhança aberta
    formal afim noetheriana $V$ tal que $f^{-1}(V)$ é
    um aberto formal afim e
    $\Gamma(f^{-1}(V),\sh{O}_{\mathfrak{X}})$ é um
    $\Gamma(V,\sh{O}_{\mathfrak{Y}})$-módulo de tipo finito.
\end{enumerate}
\end{proposition}

\begin{proof}
Segue imediatamente de \sref[I]{I.10.12.3} que \emph{a)} implica
\emph{b)}. Para ver que \emph{b)} implica \emph{c)}, podemos supor
que
\[
  \mathfrak{Y}=\Spf(B),\qquad \mathfrak{X}=\Spf(A),
\]
onde $B$ é um anel noetheriano ádico e $\mathfrak{R}$ é um ideal
de definição de $B$. Por hipótese, $X_0$ é um esquema afim
cujo anel $A_0$ é um $B/\mathfrak{R}$-módulo de tipo finito
\sref[II]{II.6.1.3}. Por \sref[I]{I.5.1.9}, cada $X_n$ é um esquema
afim; se $A_n$ denota seu anel, a condição \emph{b)} implica que,
para $m\leq n$,
\[
  A_m\simeq A_n/\mathfrak{R}^{m+1}A_n.
\]
Segue-se que $\mathfrak{X}\simeq\Spf(A)$, onde
$A=\varprojlim A_n$, e a conclusão se segue de
\sref[0]{0.7.2.9}.

Finalmente, para provar que \emph{c)} implica \emph{a)}, podemos restringir-nos
novamente ao caso
\[
  \mathfrak{Y}=\Spf(B),\qquad \mathfrak{X}=\Spf(A),
\]
onde $A$ é uma $B$-álgebra finita. Como
$A/\mathfrak{R}A$ é então uma
$B/\mathfrak{R}$-álgebra finita, \sref[I]{I.10.10.9} mostra que se
satisfazem as condições de \emph{a)}.
\end{proof}

\begin{definition}[4.8.2]
\label{III.4.8.2}
Quando se satisfazem as condições equivalentes \emph{a)}, \emph{b)} e \emph{c)}
de \sref{III.4.8.1}, diz-se que o morfismo $f$ é
\emph{finito}; também se diz que $\mathfrak{X}$ é um
\emph{$\mathfrak{Y}$-esquema formal finito}, ou um esquema
formal \emph{finito sobre $\mathfrak{Y}$}.
\end{definition}

\begin{proposition}[4.8.3]
\label{III.4.8.3}
\begin{enumerate}
  \item[\rm{(i)}] Uma imersão fechada de esquemas formais localmente
    noetherianos é um morfismo finito.
  \item[\rm{(ii)}] O composto de dois morfismos finitos de esquemas formais
    localmente noetherianos é um morfismo finito.
  \item[\rm{(iii)}] Sejam $\mathfrak{X}$, $\mathfrak{Y}$ e
    $\mathfrak{S}$ esquemas formais localmente noetherianos,
    $f:\mathfrak{X}\to\mathfrak{S}$ um morfismo finito, e
    $g:\mathfrak{Y}\to\mathfrak{S}$ um morfismo. Então a projeção
    \[
      \mathfrak{X}\times_{\mathfrak{S}}\mathfrak{Y}
      \longrightarrow\mathfrak{Y}
    \]
    é finita.
  \item[\rm{(iv)}] Seja $\mathfrak{S}$ um esquema formal localmente
    noetheriano, e sejam $\mathfrak{X}'$ e $\mathfrak{Y}'$
    esquemas formais localmente noetherianos tais que
    $\mathfrak{X}'\times_{\mathfrak{S}}\mathfrak{Y}'$ é localmente
    noetheriano. Se $\mathfrak{X}$ e $\mathfrak{Y}$ são
    $\mathfrak{S}$-esquemas formais localmente noetherianos e
    \[
      f:\mathfrak{X}\to\mathfrak{X}',\qquad
      g:\mathfrak{Y}\to\mathfrak{Y}'
    \]
    são $\mathfrak{S}$-morfismos finitos, então
    $f\times_{\mathfrak{S}}g$ é finito.
  \item[\rm{(v)}] Sejam
    \[
      f:\mathfrak{X}\to\mathfrak{Y},\qquad
      g:\mathfrak{Y}\to\mathfrak{Z}
    \]
    morfismos de esquemas formais localmente noetherianos tais que
    $g$ é de tipo finito e separado. Se $g\circ f$ é finito,
    então $f$ é finito.
\end{enumerate}
\end{proposition}

\begin{proof}
A afirmação \emph{(i)} é imediata. Utilizando o critério \emph{a)} de
\sref{III.4.8.1}, as demais afirmações se reduzem imediatamente às
proposições correspondentes para morfismos de esquemas ordinários
\sref[II]{II.6.1.5}. Deixamos os detalhes ao leitor, seguindo
o modelo de \sref[I]{I.10.13.5}.
\end{proof}

\begin{corollary}[4.8.4]
\label{III.4.8.4}
Sob as hipóteses de \sref[I]{I.10.9.9}, se $f$ é um morfismo
finito, também o é sua extensão $\widehat{f}$ às
completamentos.
\end{corollary}

\begin{corollary}[4.8.5]
\label{III.4.8.5}
Seja $\mathfrak{X}$ um esquema formal finito sobre
$\mathfrak{Y}$, e seja
$f:\mathfrak{X}\to\mathfrak{Y}$ o morfismo estrutural. Para
todo subconjunto aberto $U\subset\mathfrak{Y}$, o esquema formal
$f^{-1}(U)$ é finito sobre $U$.
\end{corollary}

\begin{proposition}[4.8.6]
\label{III.4.8.6}
Se $f:\mathfrak{X}\to\mathfrak{Y}$ é um morfismo finito de esquemas
formais localmente noetherianos, então
$f_*(\sh{O}_{\mathfrak{X}})$ é uma
$\sh{O}_{\mathfrak{Y}}$-álgebra coerente.
\end{proposition}

\begin{proof}
Podemos considerar $f$ como o limite indutivo de um sistema indutivo
$(f_n)$ de morfismos $f_n:X_n\to Y_n$. Demonstremos que os
$f_n$ são finitos e que $f_*(\sh{O}_{\mathfrak{X}})$ é
isomorfo ao limite projetivo dos
$(f_n)_*(\sh{O}_{X_n})$; isto provará a afirmação
\sref[I]{I.10.10.5}. Basta restringir-se ao caso
\[
  \mathfrak{Y}=\Spf(B),\qquad \mathfrak{X}=\Spf(A),
\]
e observar que, se $\mathfrak{R}$ é um ideal de definição
de $B$ e $A$ é um $B$-módulo de tipo finito, então
$A/\mathfrak{R}^{n+1}A$ é um módulo de tipo finito sobre
$B/\mathfrak{R}^{n+1}B$, e $A$ é o limite projetivo dos
$A/\mathfrak{R}^{n+1}A$.
\end{proof}

Reciprocamente:

\begin{proposition}[4.8.7]
\label{III.4.8.7}
Seja $\mathfrak{Y}$ um esquema formal localmente noetheriano e
$\sh{A}$ uma $\sh{O}_{\mathfrak{Y}}$-álgebra coerente. Existe
um esquema formal $\mathfrak{X}$ finito sobre $\mathfrak{Y}$,
único salvo $\mathfrak{Y}$-isomorfismo único, tal que
\[
  f_*(\sh{O}_{\mathfrak{X}})=\sh{A},
\]
onde $f:\mathfrak{X}\to\mathfrak{Y}$ é o morfismo estrutural.
\end{proposition}

\begin{proof}
\oldpage[III]{148}
Seja $\sh{K}$ um ideal de definição de $\mathfrak{Y}$, e ponhamos
\[
  Y_n=(\mathfrak{Y},
       \sh{O}_{\mathfrak{Y}}/\sh{K}^{n+1}),
  \qquad
  \sh{A}_n=\sh{A}/\sh{K}^{n+1}\sh{A}.
\]
Claramente, $\sh{A}_n$ é uma $\sh{O}_{Y_n}$-álgebra finita e, portanto,
define um $Y_n$-esquema finito
\[
  X_n=\Spec(\sh{A}_n)
\]
\sref[II]{II.6.1.3}. Para $m\leq n$, o homomorfismo canônico
sobrejetivo
\[
  h_{mn}:\sh{A}_n\longrightarrow\sh{A}_m
\]
define um morfismo $u_{mn}:X_m\to X_n$ para o qual o diagrama
\[
  \xymatrix{
    X_n\ar[d]_{f_n} &
    X_m\ar[l]_{u_{mn}}\ar[d]^{f_m}\\
    Y_n &
    Y_m\ar[l]
  }
\]
é comutativo, onde $f_n$ denota o morfismo estrutural; além disso,
identifica $X_m$ com
$X_n\times_{Y_n}Y_m$, como se segue imediatamente de
\sref[II]{II.1.4.6}. O esquema formal $\mathfrak{X}$ obtido
como limite indutivo do sistema indutivo $(X_n)$ é, portanto,
localmente noetheriano, e seu morfismo estrutural
\[
  f:\mathfrak{X}\longrightarrow\mathfrak{Y},
\]
o limite indutivo do sistema $(f_n)$, é finito
\sref{III.4.8.1} \sref[I]{I.10.12.3.1}. Além disso, a demonstração de
\sref{III.4.8.6} mostra que
$f_*(\sh{O}_{\mathfrak{X}})$ é o limite projetivo das
$\sh{A}_n$, e é, portanto, igual a $\sh{A}$
\sref[I]{I.10.10.6}. A afirmação de unicidade se segue do
resultado mais geral seguinte.
\end{proof}

\begin{proposition}[4.8.8]
\label{III.4.8.8}
Seja $\mathfrak{Y}$ um esquema formal localmente noetheriano, e
sejam $\mathfrak{X}$ e $\mathfrak{X}'$
$\mathfrak{Y}$-esquemas formais finitos sobre $\mathfrak{Y}$, com
morfismos estruturais
\[
  f:\mathfrak{X}\to\mathfrak{Y},\qquad
  f':\mathfrak{X}'\to\mathfrak{Y}.
\]
Existe uma bijeção canônica
\[
  \Hom_{\mathfrak{Y}}(\mathfrak{X},\mathfrak{X}')
  \xrightarrow{\ \sim\ }
  \Hom_{\sh{O}_{\mathfrak{Y}}}
  \bigl(f'_*(\sh{O}_{\mathfrak{X}'}),
        f_*(\sh{O}_{\mathfrak{X}})\bigr)\footnotemark.
\]
\footnotetext{A última expressão denota o conjunto de homomorfismos de
$\sh{O}_{\mathfrak{Y}}$-álgebras
$f'_*(\sh{O}_{\mathfrak{X}'})\to f_*(\sh{O}_{\mathfrak{X}})$.}
\end{proposition}

\begin{proof}
A aplicação $h\mapsto\sh{A}(h)$ se define como em
\sref[II]{II.1.1.2}. Para provar que é bijetiva, reduzimo-nos de
imediato ao caso em que
$\mathfrak{Y}=\Spf(B)$ é um esquema formal afim noetheriano.
Então
\[
  \mathfrak{X}=\Spf(A),\qquad
  \mathfrak{X}'=\Spf(A'),
\]
onde $A$ e $A'$ são $B$-álgebras finitas, e
\[
  f_*(\sh{O}_{\mathfrak{X}})=A^\Delta,\qquad
  f'_*(\sh{O}_{\mathfrak{X}'})=A'^\Delta.
\]
A conclusão se segue da correspondência bijetiva entre
$\mathfrak{Y}$-morfismos
$\mathfrak{X}\to\mathfrak{X}'$ e homomorfismos contínuos de
$B$-álgebras $A'\to A$ \sref[I]{I.10.2.2}, junto com a
correspondência bijetiva entre homomorfismos de $B$-módulos
$A'\to A$ e homomorfismos de
$\sh{O}_{\mathfrak{Y}}$-módulos
$A'^\Delta\to A^\Delta$ \sref[I]{I.10.10.2.3}.
\end{proof}

\begin{corollary}[4.8.9]
\label{III.4.8.9}
Sob a bijeção canônica de \sref{III.4.8.8}, as imersões
fechadas $\mathfrak{X}\to\mathfrak{X}'$ correspondem a
homomorfismos sobrejetivos de $\sh{O}_{\mathfrak{Y}}$-álgebras
\[
  f'_*(\sh{O}_{\mathfrak{X}'})
  \longrightarrow
  f_*(\sh{O}_{\mathfrak{X}}).
\]
\end{corollary}

\begin{proof}
A questão é local sobre $\mathfrak{Y}$ e, portanto, reduz-se à
definição das imersões fechadas de esquemas formais localmente
noetherianos \sref[I]{I.10.14.2}.
\end{proof}

\begin{corollary}[4.8.10]
\label{III.4.8.10}
Com a notação e as hipóteses de \sref{III.4.8.1}, um morfismo
ádico $f$ é uma imersão fechada se e somente se $f_0$ é uma
imersão fechada de esquemas ordinários.
\end{corollary}

\begin{proof}
Isto se segue imediatamente de \sref{III.4.8.9} e do critério de
sobrejetividade para um homomorfismo de
$\sh{O}_{\mathfrak{Y}}$-módulos coerentes \sref[I]{I.10.11.5}.
\end{proof}

\begin{proposition}[4.8.11]
\label{III.4.8.11}
Um morfismo $f:\mathfrak{X}\to\mathfrak{Y}$ de esquemas formais
localmente noetherianos
\oldpage[III]{149}
é finito se e somente se é próprio e todas suas fibras
$f^{-1}(y)$, para $y\in\mathfrak{Y}$, são finitas.
\end{proposition}

\begin{proof}
Pelas definições \sref{III.3.4.1} e \sref{III.4.8.2}, a
afirmação reduz-se imediatamente à mesma afirmação para $f_0$
(com a notação de \sref{III.4.8.1}), que é precisamente
\sref{III.4.4.2}.
\end{proof}
